\chapter{Abelian-variety rigidity: morphisms from a genus-\(0\) curve into an abelian variety are constant}
\label{chap:AbelianVarietyRigidity}
\label{chap:avr_for_rr}
% archon:covers AlgebraicJacobian/AbelianVarietyRigidity.lean AlgebraicJacobian/Genus0BaseObjects.lean AlgebraicJacobian/Genus0BaseObjects/BareScheme.lean AlgebraicJacobian/Genus0BaseObjects/ChartIso.lean AlgebraicJacobian/Genus0BaseObjects/Points.lean AlgebraicJacobian/Genus0BaseObjects/GmScaling.lean AlgebraicJacobian/RigidityLemma.lean
% Local macro: diagrammatic (Lean `≫`) composition. Defined locally because this
% chapter is the writer's only editable file; flagged in the writer report for promotion
% to macros/common.tex if desired. \providecommand is a no-op if it is later added globally.
\providecommand{\fatsemi}{\mathbin{\,;\,}}
% STRATEGY NOTE (iter-157). This chapter is the dedicated home of the project's COMMITTED
% genus-0 route (route (c)): the characteristic-free abelian-variety rigidity stack. It is
% 1:1 with the new upstream Lean file AlgebraicJacobian/AbelianVarietyRigidity.lean, which
% imports AlgebraicJacobian.Genus and is imported by AlgebraicJacobian.Jacobian (breaking the
% RigidityKbar -> Rigidity -> Jacobian import cycle so that genusZeroWitness can consume the
% keystone directly). The headline prop:rigidity_genus0_curve_to_AV is the char-free
% replacement (no [CharZero]) for the over-k-bar declaration rigidity_over_kbar of
% \texttt{chap:RigidityKbar} (which remains in tree as the fallback route (a) artifact and
% still carries [CharZero]). The df=0 / H^0(P^1,Omega)=0 / Serre-duality framing is the
% fallback (a), NOT the route developed here.
The genus-\(0\) arm of the Jacobian-witness construction has a \emph{trivial} witness object
(\(J = \Spec k\)); its only remaining mathematical content is a rigidity \emph{statement}:
every pointed morphism from a smooth proper geometrically irreducible curve of genus \(0\)
into an abelian variety is the constant morphism at the identity. We prove this
\textbf{characteristic-free} --- the protected goal is over an arbitrary field \(k\), so the
cheaper differential argument (via \(H^0(\mathbb P^1, \Omega) = H^0(\mathbb P^1, \mathcal
O(-2)) = 0\), hence \(df = 0\), hence \(f\) constant) is unavailable: it breaks on a
characteristic-\(p\) Frobenius-descent wall, where a nonconstant \(f\) can have \(df = 0\). The
rigidity route below sidesteps that wall.
The chain has two parts. The \emph{Rigidity Lemma} (\cref{thm:rigidity_lemma}, Mumford,
Form~I; equivalently Milne's Rigidity Theorem~1.1) is an elementary properness/closed-map
argument --- it uses neither the theorem of the cube nor any cohomology, is the entry point
for formalisation, and is now proven axiom-clean. From the Rigidity Lemma, Milne's \S I.3
derives the genus-\(0\) base case \emph{without} the theorem of the cube --- and, in this
project's reading, \emph{without} Milne's Theorem~3.2 (rational-map extension), Lemma~3.3,
Auslander--Buchsbaum, or differentials (the ``4th route''; see
\cref{rmk:base_case_fourth_route}). The decisive observation is that \(\mathbb P^1\) is
\emph{already complete} and \(f \colon \mathbb P^1 \to A\) is \emph{already a total morphism}, so
there is no rational defect map to extend over a complete surface and Theorem~3.2 never enters.
Instead the base case is closed by the \textbf{\(\mathbb G_m\)-scaling shortcut} (the project's
primary route, an Archon-original realisation of Milne Proposition~3.9): one uses the
\emph{multiplicative scaling} action \(\sigma_\times \colon \mathbb P^1 \times \mathbb G_m \to
\mathbb P^1\), \((x, \lambda) \mapsto \lambda x\) (the M\"obius scaling fixing both \(0\) and
\(\infty\)), which is a \emph{genuine total} scheme morphism on all of \(\mathbb P^1 \times \mathbb
G_m\) (\texttt{def:gaTranslationP1}). Normalising \(f(0) = \eta_A\) and composing \(h := \sigma_\times
\fatsemi f\), one applies additivity of \(\mathrm{Hom}(-, A)\) over a product
(\texttt{lem:hom_additivity_over_product}, Milne Corollary~1.5) --- whose Lean signature requires
only the \emph{first} factor complete --- with \(V = \mathbb P^1\) proper, \(W = \mathbb G_m\), and
base points \(0 \in \mathbb P^1\), \(1 \in \mathbb G_m\). Because \(0\) is a \emph{fixed point} of
scaling, the \(W\)-axis restriction \(\lambda \mapsto h(0, \lambda) = f(\lambda \cdot 0) = f(0)\) is
\emph{constant} at \(\eta_A\), so the decomposition collapses to \(h = p \fatsemi f\), i.e.\
\(f(\lambda x) = f(x)\). Specialising \(x = 1\) gives that \(f|_{\mathbb G_m}\) is constant, and since
\(\mathbb G_m\) is dense in \(\mathbb P^1\) with \(A\) separated, \(f\) is constant
(\texttt{prop:morphism_P1_to_AV_constant}, via the separated-target/dominant-source handle
\texttt{thm:GrpObj_eq_of_eqOnOpen}). This route uses \emph{only} the proven Corollary~1.5, a
scaling fixed point, and density --- it never forms the additive defect map, never invokes
\(\mathrm{Hom}(\mathbb G_a, A) = 0\), and so sidesteps the ``image of a group homomorphism is a
closed subgroup'' input that has no Mathlib support. The classical \(\mathbb G_a\)-additive
argument (\texttt{lem:hom_from_Ga_trivial} via the \(\mathbb G_a/\mathbb G_m\) incompatibility
\texttt{lem:hom_Ga_to_av_trivial}, the heart of Milne Proposition~3.9) is retained only as the
Milne-faithful alternative and is \emph{off} the genus-\(0\) critical path. The corollary that a
pointed regular map of abelian
varieties is a homomorphism (\texttt{lem:av_regular_map_is_hom}, Milne Corollary~1.2) is recorded
for completeness and for Route~A but is \emph{not} needed on the genus-\(0\) path; Milne's
Theorem~3.2 (\cref{thm:rational_map_to_av_extends}) is likewise demoted to a Route-A-only item
(the Albanese universal property), no longer on the genus-\(0\) critical path
(\cref{rmk:thm32_codim1_mathlib_gap}). Combining these
gives that every morphism \(\mathbb P^1 \to A\) is constant
(\texttt{prop:morphism_P1_to_AV_constant}); composing with the classification of genus-\(0\)
curves (\texttt{thm:genus_zero_curve_iso_p1}, \(C \cong \mathbb P^1\) over \(\bar k\)) yields the
headline \texttt{prop:rigidity_genus0_curve_to_AV}. The theorem of the cube (Milne \S I.5) is
\emph{not} used anywhere on this path (see \cref{rmk:cube_not_needed}); the chain invokes
neither Serre duality nor the FGA \(\Pic\) representability engine, so the genus-\(0\) arm is
fully decoupled from Route~A.
\medskip
Throughout, \(\bar k\) is an algebraically closed field of arbitrary characteristic. An
\emph{abelian variety} over \(\bar k\) is a complete (proper) irreducible group variety; in the
project's Lean encoding \(A\) is an object of \(\mathrm{Over}\,(\Spec \bar k)\) carrying
\(\mathtt{GrpObj}\), \(\mathtt{IsProper}\), \(\mathtt{Smooth}\), and \(\mathtt{GeometricallyIrreducible}\).
\section{The Rigidity Lemma}
\begin{lemma}
\leanok
[Rigidity Lemma (Mumford, Form I)]
  \label{thm:rigidity_lemma}
  \uses{lem:rigidity_eqOn_dense_open, lem:rigidity_core, lem:rigidity_snd_lift}
  
  \lean{AlgebraicGeometry.rigidity_lemma}
  % SOURCE: [Mumford, Abelian Varieties], Rigidity Lemma (Form I.), Ch.~II \S4, book p.~43 (read from references/mumford-abelian-varieties.pdf, PDF page 54)
  % SOURCE QUOTE: "RIGIDITY LEMMA. (Form I.) Let \(X\) be a complete variety, \(Y\) and \(Z\)
  % any varieties, and \(f \colon X \times Y \to Z\) a morphism such that for some
  % \(y_0 \in Y\), \(f(X \times \{y_0\})\) is a single point \(z_0\) of \(Z\). Then there is a
  % morphism \(g \colon Y \to Z\) such that if \(p_2 \colon X \times Y \to Y\) is the
  % projection, \(f = g \circ p_2\)."
  \textit{Source: Mumford, Abelian Varieties, Ch.~II \S4, Rigidity Lemma (Form~I), p.~43.}
  Let \(X\) be a complete variety, and let \(Y\), \(Z\) be any varieties over \(\bar k\). Let
  \(f \colon X \times Y \to Z\) be a morphism such that for some \(y_0 \in Y\) the image
  \(f(X \times \{y_0\})\) is a single point \(z_0 \in Z\). Then there is a morphism
  \(g \colon Y \to Z\) with \(f = g \circ p_2\), where \(p_2 \colon X \times Y \to Y\) is the
  projection; that is, \(f\) depends only on the \(Y\)-coordinate, and \(g(y) = f(x_0, y)\) for any
  fixed \(x_0 \in X\).
\end{lemma}
% SOURCE QUOTE PROOF: "PROOF. Choose any point \(x_0 \in X\), and define \(g \colon Y \to Z\) by
% \(g(y) = f(x_0, y)\). Since \(X \times Y\) is a variety, to show that \(f = g \circ p_2\), it is
% sufficient to show that these morphisms coincide on some open subset of \(X \times Y\). Let
% \(U\) be an affine open neighbourhood of \(z_0\) in \(Z\), \(F = Z - U\), and \(G = p_2(f^{-1}(F))\);
% then \(G\) is closed in \(Y\) since \(X\) is complete and hence \(p_2\) is a closed map. Further
% \(y_0 \notin G\) since \(f(X \times \{y_0\}) = \{z_0\}\). Therefore \(Y - G = V\) is a non-empty
% open subset of \(Y\). For each \(y \in V\), the complete variety \(X \times \{y\}\) gets mapped by
% \(f\) into the affine variety \(U\), and hence to a single point of \(U\). But this means that for
% any \(x \in X\), \(y \in V\), \(f(x, y) = f(x_0, y) = g \circ p_2(x, y)\), and this proves our
% assertion."
\begin{proof}
\leanok
% NOTE (iter-159 review): the iter-157 unsoundness REMAINS REPAIRED. The collapse hypothesis
% (_hf : f(X x {y0}) = {z0}) is threaded through all three helpers (rigidity_lemma ->
% rigidity_core -> rigidity_eqOn_dense_open) and GENUINELY CONSUMED in the body of
% rigidity_eqOn_dense_open (proves y0 \notin G, the load-bearing non-emptiness of V := Y - G).
% The counterexample f = fst is excluded (under f = fst, _hf forces X to collapse to a point).
% rigidity_eqOn_dense_open is sorry-free in its OWN body (hfib, the pullback-fibre fact over the
% k-bar-point y0, is closed). rigidity_snd_lift and snd_left_isClosedMap are axiom-clean (no sorryAx).
% NOTE (iter-160): the \uses edges run FORWARD along the true Lean dependency order
% (rigidity_lemma -> rigidity_eqOn_dense_open -> rigidity_eqOn_saturated_open_to_affine), so the
% not-proven status of the residual sorry propagates up and de-launders BOTH thm:rigidity_lemma and
% lem:rigidity_eqOn_dense_open. No backward edge from the leaf lemmas back up the chain.
% NOTE (iter-161): bridge 2 (rigidity_eqOn_saturated_open_to_affine) was DECOMPOSED into two named
% top-level sub-lemmas along the route-B (cohomology-free) split:
%   - Step 2, morphism_eq_of_eqAt_closedPoints: "two morphisms agreeing at every closed point of a
%     reduced Jacobson source into a separated target are equal" -- PROVEN axiom-clean (the
%     dense-closed-points coproduct probe + ext_of_isDominant; the connective Mathlib lacks).
%   - Step 1, rigidity_eqAt_closedPoint_of_proper_into_affine: Mumford's per-closed-slice constancy
%     -- the genuinely-deep residual, still `sorry`. THIS is now the chain's lone deep residual
%     (bridge 2's body is real assembly: Step 2 proven, the JacobsonSpace U instance now derivable).
% SIGNATURE CHANGE (iter-161, plan-authorized): Step 2's globalisation needs the closed points of
% the saturated open U to be DENSE, i.e. U Jacobson, which over Spec k-bar follows from (X tensor Y)
% being LOCALLY OF FINITE TYPE. So [LocallyOfFiniteType (X tensor Y).hom] is now carried on FIVE
% chain lemmas (rigidity_lemma, rigidity_core, rigidity_eqOn_dense_open,
% rigidity_eqOn_saturated_open_to_affine, rigidity_eqAt_closedPoint_of_proper_into_affine), in
% addition to [IsAlgClosed kbar]. Free downstream: curves and abelian varieties are of finite type
% over k-bar. The earlier "[IsAlgClosed] is the only added instance" claim is hereby retired.
  Fix any point \(x_0 \in X\) and define \(g \colon Y \to Z\) by \(g(y) = f(x_0, y)\) (the
  composite of \(y \mapsto (x_0, y)\) with \(f\)). Since \(X \times Y\) is a variety (in
  particular reduced and irreducible, so separated and with no embedded structure), to prove
  the equality of morphisms \(f = g \circ p_2\) it suffices to prove they agree on a non-empty
  open subset of \(X \times Y\).
  Let \(U\) be an affine open neighbourhood of \(z_0\) in \(Z\), put \(F = Z - U\) (a closed set),
  and set \(G = p_2(f^{-1}(F))\). Because \(X\) is complete, the projection
  \(p_2 \colon X \times Y \to Y\) is a \emph{closed} map (this is the defining property of
  completeness/properness); hence \(G\), the image of the closed set \(f^{-1}(F)\), is closed in
  \(Y\). The hypothesis \(f(X \times \{y_0\}) = \{z_0\} \subseteq U\) gives
  \(f^{-1}(F) \cap (X \times \{y_0\}) = \varnothing\), so \(y_0 \notin G\). Therefore
  \(V := Y - G\) is a non-empty open subset of \(Y\).
  For each \(y \in V\) we have \(X \times \{y\} \cap f^{-1}(F) = \varnothing\), i.e.
  \(f(X \times \{y\}) \subseteq U\). Thus the complete (proper) variety \(X \times \{y\}\) is
  mapped by \(f\) into the \emph{affine} variety \(U\). A proper connected variety mapping to an
  affine variety has image a single point (a global regular function on a complete connected
  variety is constant). Hence \(f(X \times \{y\})\) is one point, necessarily
  \(f(x_0, y) = g(y)\). So for all \(x \in X\) and \(y \in V\),
  \[ f(x, y) = f(x_0, y) = g(y) = (g \circ p_2)(x, y). \]
  The two morphisms \(f\) and \(g \circ p_2\) thus agree on the non-empty open
  \(X \times V\), which is dense in the irreducible \(X \times Y\); as \(Z\) is separated they
  agree everywhere. \qedhere
\end{proof}
\begin{remark}
  \label{rmk:rigidity_lemma_cube_free}
  The proof of \cref{thm:rigidity_lemma} uses \emph{only} (i) that completeness of \(X\) makes
  \(p_2\) a closed map, and (ii) that a proper connected variety has no nonconstant map to an
  affine variety. It uses neither the theorem of the cube nor any cohomology, and is valid in
  arbitrary characteristic. This is the designated formalisation entry point: it is the
  lowest-prerequisite link of the chain. (Mumford gives this as the geometric heart of the
  ``second proof'' of commutativity of an abelian variety; Milne packages the same statement
  as the Rigidity Theorem~1.1, requiring \(V\) complete and \(V \times W\) geometrically
  irreducible.)
\end{remark}
\begin{remark}
  \label{rmk:rigidity_lemma_decomposition}
  The Lean formalisation of \cref{thm:rigidity_lemma} splits into three helper steps, mirroring
  the proof above: \texttt{rigidity\_snd\_lift} (the cartesian-monoidal algebra identifying the
  collapsed map \((x, y) \mapsto (x_0, y)\) followed by \(f\) with \(g \circ p_2\));
  \texttt{rigidity\_core} (the scheme-level gluing: two maps out of the geometrically irreducible
  reduced \(X \times Y\) into the separated \(Z\) that agree on a non-empty open subset agree
  everywhere, by separatedness); and \texttt{rigidity\_eqOn\_dense\_open}
  (\cref{lem:rigidity_eqOn_dense_open}, the genuine geometry: the construction of the non-empty
  open \(X \times V\) and the slice-constancy of \(f\) on it). The collapse hypothesis
  \(f(X \times \{y_0\}) = \{z_0\}\) must be threaded through \emph{all three} helpers, since it is
  what makes \(V\) non-empty; a decomposition that drops it from the lower helpers is unsound
  (see \cref{lem:rigidity_eqOn_dense_open}).
  \emph{Status (iter-161).} Bridge~1 (the closed-map step) is \textbf{built} as the axiom-clean
  helper \texttt{snd\_left\_isClosedMap}. The fibre-over-\(y_0\) fact is \textbf{closed} (via the
  coarse \texttt{PullbackCarrier} layer), so \texttt{rigidity\_eqOn\_dense\_open} is
  \texttt{sorry}-free in its own body. Bridge~2 (slice-constancy on the saturated open \(U\),
  \cref{lem:rigidity_eqOn_saturated_open_to_affine}) has now been \textbf{decomposed} along its
  cohomology-free route-(c) split into two named sub-lemmas: the globalisation
  \cref{lem:morphism_eq_of_eqAt_closedPoints} (Step~2, ``agree at every closed point of a reduced
  Jacobson source \(\Rightarrow\) equal'') is \textbf{proven axiom-clean}, and the per-closed-slice
  constancy \cref{lem:rigidity_eqAt_closedPoint_of_proper_into_affine} (Step~1, Mumford's slice
  argument) is the chain's \emph{single genuinely-deep residual} \texttt{sorry}. The body of
  \cref{lem:rigidity_eqOn_saturated_open_to_affine} is therefore real assembly (Step~2 over Step~1).
  % NOTE: (iter-162 review) STALE as of iter-162 — Step~1 is now PROVEN and the WHOLE
  % Rigidity-Lemma chain (thm:rigidity_lemma … lem:rigidity_eqAt_closedPoint_of_proper_into_affine,
  % via the new helper lem:isIntegral_of_retract_of_integral) is sorry-free + axiom-clean
  % ({propext, Classical.choice, Quot.sound}, re-verified). The "single genuinely-deep residual
  % sorry" wording here and in the proof blocks of lem:rigidity_eqOn_dense_open / …saturated… /
  % …eqAt_closedPoint… should be refreshed by the plan/blueprint-writer to "chain closed iter-162".
  This route requires the base field to be algebraically closed \emph{and} the source to be of finite
  type: Step~2 needs the closed points of \(U\) to be dense (the Jacobson-space property), which over
  the algebraically closed \(\bar k\) holds because \(X \times Y\) is locally of finite type.
  Accordingly the formalization carries \textbf{both} \([\mathtt{IsAlgClosed}\ \bar k]\) and
  \([\mathtt{LocallyOfFiniteType}\ (X \otimes Y).\mathrm{hom}]\) across the chain lemmas
  (\texttt{rigidity\_lemma}, \texttt{rigidity\_core}, \texttt{rigidity\_eqOn\_dense\_open},
  \texttt{rigidity\_eqOn\_saturated\_open\_to\_affine},
  \texttt{rigidity\_eqAt\_closedPoint\_of\_proper\_into\_affine}); the earlier claim that
  \([\mathtt{IsAlgClosed}\ \bar k]\) is the \emph{only} added instance is retired. Both cost nothing
  downstream: curves and abelian varieties are of finite type over the algebraically closed \(\bar k\),
  and the consumers already assume it. The relative Stein-factorisation / proper-pushforward
  \(f_*\mathcal O = \mathcal O\) framing is a genuine Mathlib cohomology gap and is \emph{deliberately
  avoided}.
\end{remark}
\begin{lemma}\leanok
[Cartesian-monoidal collapse identity]
  \label{lem:rigidity_snd_lift}
  \lean{AlgebraicGeometry.rigidity_snd_lift}
  Let \(X\), \(Y\) be objects over \(\bar k\) and fix a \(\bar k\)-point \(x_0\) of \(X\).
  Post-composing the second projection \(p_2 \colon X \times Y \to Y\) with the slice section
  \(y \mapsto (x_0, y)\) equals the ``collapse the \(X\)-coordinate onto \(x_0\)'' endomorphism
  \((x, y) \mapsto (x_0, y)\) of \(X \times Y\). This is pure cartesian-monoidal algebra (no
  geometry): the product universal property distributes \(p_2\), the \(Y\)-component simplifies by
  the identity law, and the \(X\)-component collapses by uniqueness of maps to the terminal object.
\end{lemma}
\begin{lemma}\leanok
[Bridge~1: completeness makes the projection a closed map]
  \label{lem:snd_left_isClosedMap}
  \lean{AlgebraicGeometry.snd_left_isClosedMap}
  \uses{lem:projectiveLineBar_isProper}
  Let \(X\) be complete (proper) over \(\bar k\) and \(Y\) any object over \(\bar k\). Then the
  underlying topological map of the second projection \(p_2 \colon X \times Y \to Y\) is a
  \emph{closed} map. This is Mumford's ``completeness of \(X\) makes \(p_2\) a closed map'': the
  projection is the base change of \(X \to \Spec \bar k\) along \(Y \to \Spec \bar k\), properness
  of \(X\) gives universal closedness, and universal closedness is stable under base change, so the
  base-changed projection is universally closed and hence a closed map. Characteristic-free; no
  theorem of the cube, no cohomology.
\end{lemma}
\begin{lemma}\leanok
[Scheme-level gluing core of the Rigidity Lemma]
  \label{lem:rigidity_core}
  \lean{AlgebraicGeometry.rigidity_core}
  \uses{lem:rigidity_eqOn_dense_open}
  Let \(\bar k\) be algebraically closed, \(X\) proper, \(X \times Y\) geometrically irreducible,
  reduced and locally of finite type over \(\bar k\), and \(Z\) separated. Let
  \(f \colon X \times Y \to Z\) collapse the slice \(X \times \{y_0\}\) to a single point \(z_0\).
  Then \(f\) equals its composite with the collapse endomorphism \((x, y) \mapsto (x_0, y)\), i.e.\
  \(f = \mathtt{retract} \fatsemi f\). The proof takes the non-empty open \(U\) on which \(f\) and
  \(\mathtt{retract} \fatsemi f\) agree (\cref{lem:rigidity_eqOn_dense_open}); since \(X \times Y\)
  is geometrically irreducible over the one-point base \(\Spec \bar k\), it is irreducible, so the
  non-empty open \(U\) is dense and its inclusion is dominant; the dominant-source/separated-target
  rigidity handle then promotes agreement on \(U\) to equality of the two morphisms everywhere.
\end{lemma}
\begin{lemma}
\leanok
[Dense-open agreement: the geometric heart of the Rigidity Lemma]
  \label{lem:rigidity_eqOn_dense_open}
  \lean{AlgebraicGeometry.rigidity_eqOn_dense_open}
  \uses{lem:rigidity_eqOn_saturated_open_to_affine, lem:snd_left_isClosedMap}
  % SOURCE: [Mumford, Abelian Varieties], Rigidity Lemma (Form I.), Ch.~II \S4, book p.~43 (read from references/mumford-abelian-varieties.pdf, PDF page 54)
  % SOURCE QUOTE: "RIGIDITY LEMMA. (Form I.) Let \(X\) be a complete variety, \(Y\) and \(Z\)
  % any varieties, and \(f \colon X \times Y \to Z\) a morphism such that for some
  % \(y_0 \in Y\), \(f(X \times \{y_0\})\) is a single point \(z_0\) of \(Z\). Then there is a
  % morphism \(g \colon Y \to Z\) such that if \(p_2 \colon X \times Y \to Y\) is the
  % projection, \(f = g \circ p_2\)."
  \textit{Source: Mumford, Abelian Varieties, Ch.~II \S4, Rigidity Lemma (Form~I), p.~43 (the
  dense-open construction is the body of that proof).}
  Let \(X\) be a complete variety and \(Y\), \(Z\) any varieties over \(\bar k\). Let
  \(f \colon X \times Y \to Z\) be a morphism, fix a point \(x_0 \in X\), and suppose given
  \(y_0 \in Y\), \(z_0 \in Z\) with the \emph{collapse hypothesis}
  \[ f(X \times \{y_0\}) = \{z_0\} \qquad
     \text{(in Lean: } \mathtt{lift}\,(\mathbf 1_X)\,(\mathtt{toUnit}\,X \fatsemi y_0) \fatsemi f
     = \mathtt{toUnit}\,X \fatsemi z_0\text{).} \]
  Then there is a non-empty open subset \(U \subseteq X \times Y\) on which \(f\) agrees, as a
  scheme morphism, with the collapsed map \(\mathtt{retract} \fatsemi f\), where
  \(\mathtt{retract} := \mathtt{lift}\,(\mathtt{toUnit}\,(X \times Y) \fatsemi x_0)\,
  (p_2 \colon X \times Y \to Y)\) is the ``\((x, y) \mapsto (x_0, y)\)'' endomorphism. Concretely
  \(U = X \times V\) with \(V := Y - G\), \(G = p_2(f^{-1}(Z - U_0))\) for an affine open
  neighbourhood \(U_0\) of \(z_0\), and on \(U\) one has \(f(x, y) = f(x_0, y)\).
  \medskip
  \noindent\textbf{The collapse hypothesis is load-bearing.} The non-emptiness of \(V\) is exactly
  the assertion \(y_0 \notin G\), and this holds \emph{because} \(f(X \times \{y_0\}) = \{z_0\}
  \subseteq U_0\) keeps the slice over \(y_0\) inside the affine \(U_0\), away from \(F = Z - U_0\). A
  version of this lemma \emph{without} the collapse hypothesis is false: for \(X = Y = Z\) and
  \(f = p_1\) the first projection, \(f\) and \(\mathtt{retract} \fatsemi f \colon (x, y) \mapsto x_0\)
  agree on no non-empty open (the agreement locus is contained in \(\{x = x_0\}\), which has empty
  interior in the irreducible \(X \times Y\)). Therefore the formal target \emph{must} carry the
  collapse hypothesis as an explicit antecedent; dropping it (as an earlier decomposition
  silently did) makes the deferred goal unsatisfiable.
  \medskip
  \noindent\textbf{Formalization note: algebraically closed base, and locally-of-finite-type
  source.} The prose already takes \(\bar k\) algebraically closed, but the formal statement must
  \emph{encode} two properties. First, the cohomology-free proof of the slice-constancy bridge
  (below) needs each per-closed-point residue field \(\kappa(y)\) to be \(\bar k\) itself, so that a
  proper integral slice has global sections \(\Gamma = \bar k\) and maps to a single \(\bar k\)-point of
  the affine \(U_0\); this requires \([\mathtt{IsAlgClosed}\ \bar k]\). Second, the globalisation of the
  per-slice agreement (Step~2 of the bridge, \cref{lem:morphism_eq_of_eqAt_closedPoints}) needs the
  closed points of the saturated open \(U\) to be \emph{dense} --- i.e.\ \(U\) to be a Jacobson space.
  Over the algebraically closed \(\bar k\) this holds precisely because \(X \times Y\) is \emph{locally of
  finite type} over \(\bar k\) (\texttt{LocallyOfFiniteType.jacobsonSpace} transports the Jacobson
  property of \(\Spec \bar k\), and \texttt{closure\_closedPoints} then gives density); this requires
  \([\mathtt{LocallyOfFiniteType}\ (X \otimes Y).\mathrm{hom}]\). Accordingly the formalization carries
  \emph{both} \([\mathtt{IsAlgClosed}\ \bar k]\) and
  \([\mathtt{LocallyOfFiniteType}\ (X \otimes Y).\mathrm{hom}]\) on this lemma (in addition to
  \([\mathtt{Field}\ \bar k]\)), and propagates the same two instances to \cref{thm:rigidity_lemma},
  its core helper \texttt{rigidity\_core}, and the bridge sub-lemmas
  \cref{lem:rigidity_eqOn_saturated_open_to_affine} and
  \cref{lem:rigidity_eqAt_closedPoint_of_proper_into_affine}. Both cost nothing downstream: the
  consumers \texttt{prop:morphism_P1_to_AV_constant} and \texttt{prop:rigidity_genus0_curve_to_AV} (Lean
  \texttt{rigidity\_genus0\_curve\_to\_grpScheme}) \emph{already} assume
  \([\mathtt{IsAlgClosed}\ \bar k]\) --- the base-field variable is literally named \(\mathtt{kbar}\) ---
  and curves and abelian varieties are of finite type over \(\bar k\), so the
  locally-of-finite-type instance is automatic there.
\end{lemma}
% SOURCE QUOTE PROOF: "PROOF. Choose any point \(x_0 \in X\), and define \(g \colon Y \to Z\) by
% \(g(y) = f(x_0, y)\). Since \(X \times Y\) is a variety, to show that \(f = g \circ p_2\), it is
% sufficient to show that these morphisms coincide on some open subset of \(X \times Y\). Let
% \(U\) be an affine open neighbourhood of \(z_0\) in \(Z\), \(F = Z - U\), and \(G = p_2(f^{-1}(F))\);
% then \(G\) is closed in \(Y\) since \(X\) is complete and hence \(p_2\) is a closed map. Further
% \(y_0 \notin G\) since \(f(X \times \{y_0\}) = \{z_0\}\). Therefore \(Y - G = V\) is a non-empty
% open subset of \(Y\). For each \(y \in V\), the complete variety \(X \times \{y\}\) gets mapped by
% \(f\) into the affine variety \(U\), and hence to a single point of \(U\). But this means that for
% any \(x \in X\), \(y \in V\), \(f(x, y) = f(x_0, y) = g \circ p_2(x, y)\), and this proves our
% assertion."
\begin{proof}
  \leanok
  Let \(U_0\) be an affine open neighbourhood of \(z_0\) in \(Z\), put \(F = Z - U_0\), and set
  \(G = p_2(f^{-1}(F))\). Because \(X\) is complete, the projection \(p_2 \colon X \times Y \to Y\) is a
  \emph{closed} map, so \(G\) (the image of the closed set \(f^{-1}(F)\)) is closed in \(Y\). The
  collapse hypothesis \(f(X \times \{y_0\}) = \{z_0\} \subseteq U_0\) gives
  \(f^{-1}(F) \cap (X \times \{y_0\}) = \varnothing\), hence \(y_0 \notin G\); this is the only place
  the collapse hypothesis is used, and it is essential, for \(V := Y - G\) is non-empty precisely
  because \(y_0 \in V\). For each \(y \in V\) the proper variety \(X \times \{y\}\) is mapped by \(f\)
  into the affine \(U_0\), hence to a single point, necessarily \(f(x_0, y)\). So for all \(x \in X\),
  \(y \in V\), we have \(f(x, y) = f(x_0, y) = (\mathtt{retract} \fatsemi f)(x, y)\), i.e.\ \(f\) and
  \(\mathtt{retract} \fatsemi f\) agree on the non-empty open \(U := X \times V\). (In the Lean
  encoding the two slice-constancy bridges are the closed-map step from properness and the
  ``proper-into-affine has single-point image'' step; see the docstring of
  \texttt{rigidity\_core}.)
  \medskip
  \noindent\textbf{Formalization notes (the concrete char-free Mathlib route).} The textbook
  proof above is realised in Lean by three pieces. They cite Mathlib lemma names as the chosen
  route; none uses cohomology, and the relative Stein-factorisation framing is deliberately
  \emph{avoided} (see below).
  \emph{Bridge~1 (the closed-map step) --- BUILT.} ``\(X\) complete \(\Rightarrow p_2\) closed'' is
  the axiom-clean helper \texttt{snd\_left\_isClosedMap} (iter-158). It is obtained from
  \texttt{IsProper.toUniversallyClosed} on \(X\), base-changed across the canonical pullback square
  to make \(p_2\) universally closed, using that the monoidal projection \emph{is} the chosen
  pullback projection: \texttt{Over.snd\_left} rewrites \((\mathtt{snd}\,X\,Y).\mathtt{left}\) to
  \(\mathtt{Limits.pullback.snd}\ X.\mathtt{hom}\ Y.\mathtt{hom}\) exactly. This step is already
  proven.
  \emph{The fibre fact (non-emptiness of \(V\)) --- RESOLVED route, char-free.} That the slice
  \(X \times \{y_0\}\) over the \(\bar k\)-rational point \(y_0\) is exactly the image of the slice
  section \(s \colon X \to X \times Y\) is formalised \emph{without} the residue-field /
  \texttt{Triplet} / tensor-product machinery (that fine \texttt{carrierEquiv} layer is the wrong
  granularity and is to be avoided). One pastes the identity square for the section against the
  canonical pullback square (\texttt{IsPullback.of\_right},
  \texttt{IsPullback.of\_horiz\_isIso}; the section identities \(s \fatsemi p_1 = \mathbf 1\) and
  \(s \fatsemi p_2 = X.\mathtt{hom} \fatsemi y_0\) come from \texttt{lift\_fst}/\texttt{lift\_snd}
  with \texttt{Over.w} on the unit) and reads off the fibre with the \emph{coarse} topological
  layer of \texttt{PullbackCarrier}, namely
  \texttt{AlgebraicGeometry.Scheme.Pullback.image\_preimage\_eq\_of\_isPullback}. This needs no
  \([\mathtt{IsAlgClosed}]\).
  \emph{Bridge~2 (slice-constancy / the agreement equation) --- RESOLVED route, cohomology-FREE.}
  ``Proper connected slice into affine \(\Rightarrow\) single point, glued to a morphism equality on
  \(U\)'' is realised by a \textbf{global-sections + per-closed-point} route, \emph{not} by any
  relative Stein-factorisation or proper-pushforward \(f_*\mathcal O = \mathcal O\) statement: the
  latter is a genuine Mathlib gap (no Stein factorisation, no \(H^0\) base change), and route~(c)
  exists precisely to avoid coherent cohomology, so it must \emph{not} be attempted. The concrete
  stack is: for each closed point \(y \in V\) one has \(\kappa(y) = \bar k\) (alg-closed, finite type);
  the slice \(X_y\) is proper integral, so its global sections form a field
  (\texttt{isField\_of\_universallyClosed}) module-finite over \(\bar k\)
  (\texttt{finite\_appTop\_of\_universallyClosed}), and algebraic closedness kills the finite
  extension to give \(\Gamma(X_y) = \bar k\). A morphism into the affine \(U_0\) is pinned by its ring
  map on global sections (\texttt{ext\_of\_isAffine}), so \(f\) and
  \(\mathtt{retract} \fatsemi f\) agree on each closed slice. Closed points are dense in the
  locally-of-finite-type \(U\) (\texttt{closure\_closedPoints} / the Jacobson-space property of
  locally-of-finite-type \(\bar k\)-schemes; in the formalization this density is supplied by the new
  hypothesis \([\mathtt{LocallyOfFiniteType}\ (X \otimes Y).\mathrm{hom}]\)), so the per-slice
  agreement globalises to all of \(U\) via the reduced-source / separated-target rigidity
  \texttt{ext\_of\_isDominant\_of\_isSeparated'} (the same handle \texttt{rigidity\_core} already
  uses). The packaged ``morphisms agreeing on a dense set of closed points are equal''
  hom-extensionality --- the one connective Mathlib does not supply directly (it gives only the
  single-dominant-morphism \texttt{ext\_of\_isDominant}) --- has now been built as the axiom-clean
  \cref{lem:morphism_eq_of_eqAt_closedPoints} (Step~2): it assembles the closed points into one
  dominant coproduct probe and applies \texttt{ext\_of\_isDominant}. The whole argument is extracted
  into the dedicated obligation \cref{lem:rigidity_eqOn_saturated_open_to_affine} below, whose body
  wires that proven Step~2 over the per-slice Step~1
  (\cref{lem:rigidity_eqAt_closedPoint_of_proper_into_affine}); Step~1 (Mumford's per-closed-slice
  constancy) is now the chain's \emph{single} genuinely-deep residual \texttt{sorry}.
\end{proof}
\begin{lemma}
\leanok
[Bridge~2: slice-constancy on a saturated open mapping into an affine]
  \label{lem:rigidity_eqOn_saturated_open_to_affine}

  \lean{AlgebraicGeometry.rigidity_eqOn_saturated_open_to_affine}
  \uses{lem:rigidity_eqAt_closedPoint_of_proper_into_affine, lem:morphism_eq_of_eqAt_closedPoints}
  % NOTE (iter-161): the iter-160 SIGNATURE GAP is RESOLVED. The Lean proof wires Step 2
  % (morphism_eq_of_eqAt_closedPoints, PROVEN axiom-clean) over Step 1
  % (rigidity_eqAt_closedPoint_of_proper_into_affine, the residual sorry). Step 2 needs the closed
  % points of U to be DENSE, i.e. U Jacobson; over k̄ this follows from (X⊗Y) being locally of finite
  % type, so [LocallyOfFiniteType (X⊗Y).hom] is now a stated hypothesis of this lemma (and across the
  % chain). With it, the `JacobsonSpace U` instance is a routine discharge (LocallyOfFiniteType.
  % jacobsonSpace + JacobsonSpace.of_isOpenEmbedding), no longer an as-typed gap. The only `sorry`
  % reachable from this lemma is now Step 1 (Mumford's per-closed-slice constancy). The hypothesis
  % set is recorded in the statement and proof prose below (both [IsAlgClosed k̄] AND
  % [LocallyOfFiniteType (X⊗Y).hom]); the earlier "[IsAlgClosed] is the only added instance" claim
  % is retired.
  % SOURCE: [Mumford, Abelian Varieties], Rigidity Lemma (Form I.), Ch.~II \S4, book p.~43 (read from references/mumford-abelian-varieties.pdf, PDF page 54)
  % SOURCE QUOTE PROOF: "For each \(y \in V\), the complete variety \(X \times \{y\}\) gets mapped by
  % \(f\) into the affine variety \(U\), and hence to a single point of \(U\). But this means that for
  % any \(x \in X\), \(y \in V\), \(f(x, y) = f(x_0, y) = g \circ p_2(x, y)\), and this proves our
  % assertion." [Mumford, Ch.~II \S4, Rigidity Lemma, p.~43.]
  \textit{Source: Mumford, Abelian Varieties, Ch.~II \S4, Rigidity Lemma (Form~I), p.~43 (the
  ``for each \(y \in V\), the complete slice maps into the affine, hence to a single point'' step).}
  Let \(\bar k\) be algebraically closed and let \(X\), \(Y\), \(Z\) be objects of
  \(\mathrm{Over}\,(\Spec \bar k)\) with \(X\) proper, \(X \times Y\) geometrically irreducible, reduced,
  and \emph{locally of finite type} over \(\bar k\) (the formalization carries
  \([\mathtt{LocallyOfFiniteType}\ (X \otimes Y).\mathrm{hom}]\), which makes the saturated open \(U\)
  below a Jacobson space, so its closed points are dense --- the hypothesis Step~2 of the proof
  consumes), and \(Z\) separated. Let \(f \colon X \times Y \to Z\) be a morphism, fix a \(\bar k\)-point
  \(x_0 \colon \mathbf 1 \to X\), and write \(\mathtt{retract} := \mathtt{lift}\,
  (\mathtt{toUnit}\,(X \times Y) \fatsemi x_0)\,(p_2)\) for the ``\((x, y) \mapsto (x_0, y)\)''
  endomorphism. Suppose:
  \begin{itemize}
    \item \(U \subseteq X \times Y\) is a \(p_2\)-\emph{saturated} open, i.e.\ \(U = p_2^{-1}(V_{\mathrm{set}})\)
      for a set \(V_{\mathrm{set}} \subseteq Y\) (so \(U\) contains the whole fibre \(X_y\) over each of
      its points);
    \item \(U_0 \subseteq Z\) is an \emph{affine} open;
    \item \(f(U) \subseteq U_0\), i.e.\ \(f\) maps \(U\) into the affine \(U_0\).
  \end{itemize}
  Then \(f\) and \(\mathtt{retract} \fatsemi f\) agree on \(U\) as scheme morphisms:
  \(U.\iota \fatsemi f = U.\iota \fatsemi (\mathtt{retract} \fatsemi f)\).
\end{lemma}
\begin{proof}
  \leanok
  This is the formalization realization of Mumford's step ``for each \(y \in V\), the complete
  variety \(X \times \{y\}\) gets mapped by \(f\) into the affine variety \(U\), and hence to a single
  point of \(U\)''. We realise it by a \textbf{global-sections + per-closed-point} route that is
  \emph{cohomology-free}; in particular the relative Stein-factorisation / proper-pushforward
  \(f_*\mathcal O = \mathcal O\) framing is a confirmed Mathlib gap (no Stein factorisation, no
  \(H^0\) base change) and is \emph{deliberately avoided}. The argument splits into two named
  sub-lemmas, one residual and one proven.
  \medskip
  \noindent\textbf{Step 1 (per-closed-slice constancy) --- \cref{lem:rigidity_eqAt_closedPoint_of_proper_into_affine}, the residual.}
  Fix a closed point \(x\) of \(U\), lying over a closed point \(y = p_2(x) \in V_{\mathrm{set}}\).
  Since \(\bar k\) is algebraically closed and \(Y\) is of finite type over \(\bar k\), the residue field
  is \(\kappa(y) = \bar k\). Because \(U = p_2^{-1}(V_{\mathrm{set}})\) is saturated, \(U\) contains the
  \emph{entire} fibre \(X_y \cong X\) over \(y\); hence \(f\) maps the proper integral slice \(X_y\) into the
  affine \(U_0\). By \(\mathtt{isField\_of\_universallyClosed}\) the global sections \(\Gamma(X_y)\) form a
  field, module-finite over \(\bar k\) by \(\mathtt{finite\_appTop\_of\_universallyClosed}\), and
  algebraic closedness collapses the finite extension to give \(\Gamma(X_y) = \bar k\). A morphism into
  the affine \(U_0\) is pinned by its ring map on global sections (\(\mathtt{ext\_of\_isAffine}\)), so the
  restriction of \(f\) to \(X_y\) factors through a single \(\bar k\)-point of \(U_0\) --- necessarily the
  point \(f(x_0, y) = (\mathtt{retract} \fatsemi f)(x)\). Thus \(f\) and \(\mathtt{retract} \fatsemi f\)
  agree (after the residue-field probe) at every closed point of \(U\). This per-slice statement is the
  chain's single genuinely-deep residual \texttt{sorry}.
  \medskip
  \noindent\textbf{Step 2 (globalisation over dense closed points) --- \cref{lem:morphism_eq_of_eqAt_closedPoints}, proven.}
  The closed points are dense in the locally-of-finite-type \(\bar k\)-scheme \(U\)
  (\(\mathtt{closure\_closedPoints}\) / the Jacobson-space property, supplied here by the new hypothesis
  \([\mathtt{LocallyOfFiniteType}\ (X \otimes Y).\mathrm{hom}]\)). To turn ``\(f\) and
  \(\mathtt{retract} \fatsemi f\) agree at each closed point'' into a single morphism equality on \(U\),
  the proven globalisation \cref{lem:morphism_eq_of_eqAt_closedPoints} assembles the closed points
  into one dominant probe --- the coproduct
  \(\coprod_{x \in \mathrm{closedPoints}\,U} \Spec \kappa(x) \to U\), whose range is dense --- and feeds
  it to the reduced-source / separated-target rigidity handle \(\mathtt{ext\_of\_isDominant}\) (the same
  family \(\mathtt{rigidity\_core}\) uses). This yields the morphism equality
  \(U.\iota \fatsemi f = U.\iota \fatsemi (\mathtt{retract} \fatsemi f)\) on all of \(U\). The
  ``agrees on a dense set of closed points \(\Rightarrow\) hom-equality'' connective is the one piece
  Mathlib does not package directly (it supplies only the single-dominant-morphism
  \(\mathtt{ext\_of\_isDominant}\)); that bespoke build is now done in
  \cref{lem:morphism_eq_of_eqAt_closedPoints}, and it is \emph{not} cohomology. \qedhere
\end{proof}
\begin{lemma}\leanok
[Step~2: morphisms agreeing at every closed point of a Jacobson source are equal]
  \label{lem:morphism_eq_of_eqAt_closedPoints}
  
  \lean{AlgebraicGeometry.morphism_eq_of_eqAt_closedPoints}
  % Project-bespoke assembly: NO external source. Mathlib supplies only the single-dominant-morphism
  % ext_of_isDominant; the "dense closed points => hom-extensionality" connective is built here.
  Let \(W\), \(Z\) be schemes with \(W\) \emph{reduced}, \(W\) a \emph{Jacobson space}, and \(Z\)
  \emph{separated}. Let \(g_1, g_2 \colon W \to Z\) be morphisms. If, for every \emph{closed} point
  \(x \in W\), the two morphisms agree after restriction along the residue-field point
  \(\Spec \kappa(x) \to W\) (in Lean, \(W.\mathtt{fromSpecResidueField}\,x \fatsemi g_1 =
  W.\mathtt{fromSpecResidueField}\,x \fatsemi g_2\)), then \(g_1 = g_2\).
\end{lemma}
\begin{proof}
\leanok
  This is the ``dense closed points \(\Rightarrow\) hom-extensionality'' connective that Mathlib does
  not package directly: it supplies only the single-dominant-morphism extensionality
  \(\mathtt{ext\_of\_isDominant}\) (a dominant morphism into a separated target is an epimorphism for
  the purpose of right-cancellation). We assemble the per-closed-point hypotheses into one such
  dominant probe.
  Form the coproduct \(P := \coprod_{x \in \mathrm{closedPoints}\,W} \Spec \kappa(x)\) and the morphism
  \(\mathrm{probe} \colon P \to W\) defined by \(\mathtt{Sigma.desc}\) of the residue-field points
  \(W.\mathtt{fromSpecResidueField}\,x\). Its topological range contains every closed point of \(W\):
  the inclusion \(\iota_x \colon \Spec \kappa(x) \to P\) composes with \(\mathrm{probe}\) to
  \(W.\mathtt{fromSpecResidueField}\,x\) (the defining property of \(\mathtt{Sigma.desc}\)), and the
  range of \(W.\mathtt{fromSpecResidueField}\,x\) is the singleton \(\{x\}\)
  (\(\mathtt{Scheme.range\_fromSpecResidueField}\)). Because \(W\) is a Jacobson space, the closed points
  are dense (\(\mathtt{closure\_closedPoints}\), i.e.\ \(\overline{\mathrm{closedPoints}\,W} = W\)), so
  the range of \(\mathrm{probe}\) is dense and \(\mathrm{probe}\) is \emph{dominant}.
  Componentwise, \(g_1\) and \(g_2\) agree after \(\mathrm{probe}\): precomposing with each coproduct
  inclusion \(\iota_x\) reduces (by \(\mathtt{Sigma.\iota\_desc}\)) to
  \(W.\mathtt{fromSpecResidueField}\,x \fatsemi g_1 = W.\mathtt{fromSpecResidueField}\,x \fatsemi g_2\),
  which is the hypothesis at the closed point \(x\). By the universal property of the coproduct
  (\(\mathtt{Sigma.hom\_ext}\)) this gives \(\mathrm{probe} \fatsemi g_1 = \mathrm{probe} \fatsemi g_2\).
  Since \(\mathrm{probe}\) is dominant and \(Z\) is separated, \(\mathtt{ext\_of\_isDominant}\) cancels it
  on the left to yield \(g_1 = g_2\).
\end{proof}
\begin{lemma}
\leanok
[Global-sections constancy: two \(\bar k\)-points of a proper integral source into an affine agree]
  \label{lem:eq_comp_of_isAffine_of_properIntegral}
  \lean{AlgebraicGeometry.eq_comp_of_isAffine_of_properIntegral}
  % Project-bespoke assembly of Mathlib lemmas: NO external source. Realises, cohomology-free, the
  % classical "a global regular function on a proper integral variety over an algebraically closed
  % field is constant" --- precisely the affine-image step of Mumford's Rigidity Lemma proof quoted
  % above (the "complete slice into an affine maps to a single point" sentence, Ch.~II \S4, p.~43).
  Let \(\bar k\) be algebraically closed and let \(W\) be an integral scheme equipped with a structure
  morphism \(w \colon W \to \Spec \bar k\) that is universally closed (proper) and locally of finite
  type. Let \(g \colon W \to V\) be a morphism into an \emph{affine} scheme \(V\). Then any two
  \(\bar k\)-points of \(W\) --- i.e.\ any two sections \(a, b \colon \Spec \bar k \to W\) of \(w\) (so
  \(a \fatsemi w = \mathbf 1\) and \(b \fatsemi w = \mathbf 1\)) --- satisfy \(a \fatsemi g = b \fatsemi g\).
  \medskip
  \noindent\textbf{Every hypothesis is load-bearing} (confirmed by the iter-161 audit).
  \([\mathtt{IsAlgClosed}\ \bar k]\) collapses the finite extension \(\Gamma(W)/\bar k\); integrality of
  \(W\) makes \(\Gamma(W)\) a domain (without it a two-point disjoint union is a counterexample);
  \(\mathtt{UniversallyClosed}\) is what upgrades \(\Gamma(W)\) to a \emph{field} (without it
  \(\mathbb A^1\), with \(\Gamma = \bar k[t]\) and two distinct \(\bar k\)-points, is a counterexample);
  \(\mathtt{LocallyOfFiniteType}\) gives finiteness of \(\Gamma(W)\) over \(\bar k\); and
  \([\mathtt{IsAffine}\ V]\) is what lets \(\mathtt{ext\_of\_isAffine}\) pin the morphism by its ring map
  on global sections.
\end{lemma}
\begin{proof}
\leanok
  Because \(W\) is integral and \(w\) is universally closed, the global sections
  \(\Gamma(W, \mathcal O_W)\) form a field (\(\mathtt{isField\_of\_universallyClosed}\)), and the structure
  ring map \(w^\sharp \colon \bar k \to \Gamma(W)\) (the map \(\mathtt{wk.appTop}\) after the canonical
  identification \(\Gamma(\Spec \bar k) \cong \bar k\)) is integral / module-finite
  (\(\mathtt{finite\_appTop\_of\_universallyClosed}\)). Since \(\bar k\) is algebraically closed and
  \(\Gamma(W)\) is a domain finite over it, that map is bijective
  (\(\mathtt{IsAlgClosed.ringHom\_bijective\_of\_isIntegral}\)); thus \(\Gamma(W) = \bar k\) and
  \(w^\sharp\) is an isomorphism on global sections --- in particular an epimorphism. Both sections
  \(a\), \(b\) left-invert \(w^\sharp\) on global sections (they are sections of \(w\)), so they induce the
  \emph{same} ring map \(a^\sharp = b^\sharp\) on \(\Gamma(W)\) (cancel the epimorphism \(w^\sharp\)). A
  morphism into the \emph{affine} \(V\) is pinned by its global-sections ring map
  (\(\mathtt{ext\_of\_isAffine}\)), whence \(a \fatsemi g = b \fatsemi g\). This realises, without
  cohomology, the classical ``proper connected variety into an affine is a single point'' step
  quoted from Mumford above.
\end{proof}
\begin{lemma}
\leanok
[Integrality descends to a retract of an integral scheme]
  \label{lem:isIntegral_of_retract_of_integral}
  \lean{AlgebraicGeometry.isIntegral_of_retract}
  % NOTE: (iter-162 review) project-bespoke (elementary topology + algebra); NO external source.
  % LANDED axiom-clean as `AlgebraicGeometry.isIntegral_of_retract`, stated generically for any
  % retract `r : S ⟶ T`, `pr : T ⟶ S` with `r ≫ pr = 𝟙 S` of an integral `T` (more general than
  % the `X`/`X×Y`/`p_1` instance below). PROOF DIVERGENCE (cosmetic, sound either way): the Lean
  % proves reducedness PER-STALK (each `S.stalk x` embeds via `pr.stalkMap` into the reduced
  % `T.stalk`, then `isReduced_of_isReduced_stalk`), not via the global-sections split-injection
  % the prose below describes. Plan/blueprint-writer may align the prose to the stalk route.
  Let \(X \times Y\) be an integral scheme (here it is geometrically irreducible and reduced over
  \(\bar k\)). Suppose \(X\) is a \emph{retract} of \(X \times Y\): there is a section
  \(s \colon X \to X \times Y\) of the first projection \(p_1\), i.e.\ \(s \fatsemi p_1 = \mathbf 1_X\).
  Then \(X\) is integral.
\end{lemma}
\begin{proof}
\leanok
  Integrality is reducedness together with irreducibility, and the retract supplies each.
  \emph{Reduced.} The section identity \(s \fatsemi p_1 = \mathbf 1_X\) makes the pullback on global
  sections \(p_1^\sharp \colon \mathcal O_X(X) \to \mathcal O_{X \times Y}(X \times Y)\) split
  injective (it is left-inverted by \(s^\sharp\)). A split-injective ring map into the reduced ring
  \(\mathcal O_{X \times Y}(X \times Y)\) embeds \(\mathcal O_X(X)\) into a reduced ring, so it too is
  reduced (\(\mathtt{isReduced\_of\_injective}\) on the split injection); hence \(X\) is reduced.
  \emph{Irreducible.} The first projection \(p_1 \colon X \times Y \to X\) is surjective because it
  admits the section \(s\) (every point of \(X\) is \(p_1(s(\cdot))\)). The continuous image of an
  irreducible space under a surjection is irreducible; as \(X \times Y\) is irreducible and \(p_1\) is a
  continuous surjection onto \(X\), the space of \(X\) is irreducible.
  Reduced and irreducible together give that \(X\) is integral. This is elementary --- no cohomology,
  no external citation --- and feeds the Step-1 geometric assembly below, where the proper slice
  \(X_y \cong X\) must be presented as proper \emph{integral}.
\end{proof}
\begin{lemma}\leanok
[Step~1: per-closed-slice constancy of a proper slice into an affine]
  \label{lem:rigidity_eqAt_closedPoint_of_proper_into_affine}

  \lean{AlgebraicGeometry.rigidity_eqAt_closedPoint_of_proper_into_affine}
  \uses{lem:eq_comp_of_isAffine_of_properIntegral, lem:isIntegral_of_retract_of_integral}
  % SOURCE: [Mumford, Abelian Varieties], Rigidity Lemma (Form I.), Ch.~II \S4, book p.~43 (read from references/mumford-abelian-varieties.pdf, PDF page 54)
  % SOURCE QUOTE PROOF: "For each \(y \in V\), the complete variety \(X \times \{y\}\) gets mapped by
  % \(f\) into the affine variety \(U\), and hence to a single point of \(U\). But this means that for
  % any \(x \in X\), \(y \in V\), \(f(x, y) = f(x_0, y) = g \circ p_2(x, y)\), and this proves our
  % assertion." [Mumford, Ch.~II \S4, Rigidity Lemma, p.~43.]
  \textit{Source: Mumford, Abelian Varieties, Ch.~II \S4, Rigidity Lemma (Form~I), p.~43 (the
  ``for each \(y \in V\), the complete slice maps into the affine, hence to a single point'' step).}
  With the data of \cref{lem:rigidity_eqOn_saturated_open_to_affine} --- \(\bar k\) algebraically
  closed, \(X\) proper over \(\bar k\), \(X \times Y\) geometrically irreducible, reduced, and locally of
  finite type, \(Z\) separated, a \(\bar k\)-point \(x_0 \colon \mathbf 1 \to X\), a morphism
  \(f \colon X \times Y \to Z\), a \(p_2\)-saturated open \(U = p_2^{-1}(V_{\mathrm{set}})\) mapping into an
  affine open \(U_0 \subseteq Z\) --- fix a \emph{closed} point \(x\) of \(U\). Then \(f\) and the collapsed
  map \(\mathtt{retract} \fatsemi f\) (with \(\mathtt{retract} := \mathtt{lift}\,
  (\mathtt{toUnit}\,(X \times Y) \fatsemi x_0)\,(p_2)\), i.e.\ \((x, y) \mapsto (x_0, y)\)) agree at \(x\)
  after the residue-field probe \(\Spec \kappa(x) \to U\):
  \[ U.\mathtt{fromSpecResidueField}\,x \fatsemi (U.\iota \fatsemi f)
     = U.\mathtt{fromSpecResidueField}\,x \fatsemi (U.\iota \fatsemi (\mathtt{retract} \fatsemi f)). \]
\end{lemma}
\begin{proof}
  \leanok
  This is Mumford's per-slice step, realised \emph{cohomology-free} (route~B): the relative Stein /
  proper-pushforward \(f_*\mathcal O = \mathcal O\) framing is a confirmed Mathlib gap and is
  deliberately avoided. The closed point \(x\) of \(U\) lies over a closed point
  \(y = p_2(x) \in V_{\mathrm{set}}\); since \(\bar k\) is algebraically closed and \(Y\) is of finite type
  over \(\bar k\), the residue field is \(\kappa(y) = \bar k\). Saturation \(U = p_2^{-1}(V_{\mathrm{set}})\)
  puts the \emph{entire} proper integral fibre \(X_y \cong X\) over \(y\) inside \(U\), so \(f\) maps \(X_y\)
  into the affine \(U_0\). Because \(X_y\) is proper integral over \(\bar k\), its global sections form a
  field (\(\mathtt{isField\_of\_universallyClosed}\)), module-finite over \(\bar k\)
  (\(\mathtt{finite\_appTop\_of\_universallyClosed}\)); algebraic closedness collapses the finite
  extension, giving \(\Gamma(X_y, \mathcal O) = \bar k\).
  % NOTE (iter-162): the iter-161 documentary gap is RESOLVED --- this algebraic core ("two k-bar
  % points of a proper integral l.f.t. k-bar-scheme into an affine agree") now has its own node
  % lem:eq_comp_of_isAffine_of_properIntegral (above), and this proof's  edge points at it.
  A morphism into the affine \(U_0\) is pinned by
  its ring map on global sections (\(\mathtt{ext\_of\_isAffine}\)), so \(f|_{X_y}\) factors through a
  single \(\bar k\)-point of \(U_0\). That point is \(f(x_0, y)\), since the \(\bar k\)-point \((x_0, y)\) lies
  on \(X_y\); and \(f(x_0, y) = (\mathtt{retract} \fatsemi f)(x)\) by the definition of \(\mathtt{retract}\).
  Evaluating at the residue-field probe of the closed point \(x\) gives the asserted equality. This
  per-slice statement is the chain's single genuinely-deep residual \texttt{sorry}; its globalisation
  over the dense closed points is \cref{lem:morphism_eq_of_eqAt_closedPoints}.
\end{proof}
\section{The Milne \S I.3 chain: additivity, homomorphisms, extension}
This section assembles, from the proven Rigidity Lemma, the genus-\(0\) base case. The two
load-bearing ingredients of the \textbf{primary route} are the additive decomposition over a
product (\texttt{lem:hom_additivity_over_product}, Milne Corollary~1.5) and the concrete \(\mathbb
G_m\)-scaling action \(\sigma_\times\) on \(\mathbb P^1\) (\texttt{def:gaTranslationP1}); from these
the base case ``every morphism \(\mathbb P^1 \to A\) is constant'' falls out \emph{in one shot}
(\texttt{prop:morphism_P1_to_AV_constant}): the scaling fixed point \(0\) collapses one axis of the
decomposition, giving \(f(\lambda x) = f(x)\) and hence constancy on the dense \(\mathbb G_m\). The
classical \(\mathbb G_a\)-additive homomorphism step
(\texttt{lem:hom_from_Ga_trivial}) and the triviality lemma \(\mathrm{Hom}(\mathbb G_a, A) = 0\)
(\texttt{lem:hom_Ga_to_av_trivial}) are retained as the Milne-faithful alternative but are
\emph{not} on the genus-\(0\) critical path. None of this uses the theorem of the cube, and --- crucially
--- none of it uses Milne's Theorem~3.2 / Lemma~3.3 (the rational-map extension): because \(f
\colon \mathbb P^1 \to A\) is \emph{already} a total morphism out of the complete \(\mathbb
P^1\), there is no defect map to extend over a complete surface (see
\cref{rmk:base_case_fourth_route}). Corollary~1.2 (\texttt{lem:av_regular_map_is_hom}) and
Theorem~3.2 (\cref{thm:rational_map_to_av_extends}) are retained in this chapter only as
Route-A inputs (the Albanese universal property) and are explicitly \emph{off} the genus-\(0\)
path. The genus-\(0\) output is exactly the content of Milne's Proposition~3.9, ``every rational
(here: regular) map \(\mathbb P^1 \to A\) is constant''.
\begin{remark}
  \label{rmk:cube_not_needed}
  \textbf{The theorem of the cube is not needed on the genus-\(0\) path.} An earlier draft of
  this chapter recorded the theorem of the cube as a deferred prerequisite of the single-curve
  base case ``\(\mathbb P^1 \to A\) constant''. That was a misreading. The base case
  (Proposition~3.9, the \(d = 1\) case of Proposition~3.10) follows directly from the Rigidity
  Theorem~1.1 --- our \cref{thm:rigidity_lemma} --- via Corollary~1.5 (additivity,
  \texttt{lem:hom_additivity_over_product}) composed with the total \(\mathbb G_m\)-scaling
  action \(\sigma_\times\) (\texttt{def:gaTranslationP1}): the scaling fixed point \(0\) collapses one
  axis of the decomposition, yielding \(f(\lambda x) = f(x)\) and hence constancy on the dense
  \(\mathbb G_m\) (the primary \(\mathbb G_m\)-scaling shortcut; see
  \cref{rmk:base_case_fourth_route}). This route does not even use the triviality of
  \(\mathrm{Hom}(\mathbb G_a, A)\) (\texttt{lem:hom_Ga_to_av_trivial}), which is retained only for the
  Milne-faithful \(\mathbb G_a\)-additive alternative. Neither the theorem of the cube nor Milne's Theorem~3.2 /
  Lemma~3.3 (rational-map extension) is used: since \(f \colon \mathbb P^1 \to A\) is already a
  total morphism out of the complete \(\mathbb P^1\), no defect map needs extending over a complete
  surface (this corrects the earlier claim that the surface extension was on the genus-\(0\) path;
  see \cref{rmk:thm32_codim1_mathlib_gap}). The theorem of the cube (Milne \S I.5) enters only for
  projectivity, the seesaw/square principle, dual abelian varieties, polarizations, and the
  Poincar\'e sheaf --- none on the genus-\(0\) path. For Route~A: the Albanese universal property
  that the positive-genus arm needs (Milne Proposition~6.1/6.4, \S III.6) is \emph{also} cube-free,
  resting on Theorem~3.2 (\cref{thm:rational_map_to_av_extends}) + Corollary~1.2/1.5 --- which is
  precisely why Theorem~3.2 and Corollary~1.2 are retained in this chapter as Route-A inputs.
\end{remark}
% SOURCE QUOTE PROOF: "PROOF. Set f=h|V×{w0}, g=h|{v0}×W, and identify V×{w0} and {v0}×W with
% V and W. On points, f(v) = h(v,w0) and g(w) = h(v0,w), and so φ =df h − (f∘p + g∘q) is the
% map that sends (v,w) ↦ h(v,w) − h(v,w0) − h(v0,w). Thus φ(V×{w0}) = 0 = φ({v0}×W) and so the
% theorem shows that φ = 0." [Milne, Cor 1.5, \S I.1, doc p.~9--10.]
\begin{proof}
  \leanok
  \textit{Source: Milne, Abelian Varieties, Corollary~1.5.}
  Define \(f := h|_{V \times \{w_0\}}\) and \(g := h|_{\{v_0\} \times W}\), identifying
  \(V \times \{w_0\} \cong V\) and \(\{v_0\} \times W \cong W\); on points \(f(v) = h(v, w_0)\) and
  \(g(w) = h(v_0, w)\). Form the difference
  \[ \varphi := h - (f \circ p + g \circ q) \colon V \times W \to A, \qquad
     \varphi(v, w) = h(v, w) - h(v, w_0) - h(v_0, w) \]
  (the subtraction is in the group \(A\), i.e.\ \(\varphi\) is the composite of
  \((h, -(f \circ p + g \circ q))\) with the group law of \(A\)). Then \(\varphi\) vanishes on both
  axes: \(\varphi(V \times \{w_0\}) = 0 = \varphi(\{v_0\} \times W)\) (direct computation, using
  \(h(v_0, w_0) = 0\)). In particular \(\varphi\) collapses the complete axis \(V \times \{w_0\}\) to
  the single point \(0 \in A\), so by the Rigidity Lemma (\cref{thm:rigidity_lemma}, with \(V\)
  complete and \(V \times W\) geometrically irreducible) \(\varphi\) factors through the projection
  \(q\); but it also vanishes on \(\{v_0\} \times W\), which is a section of \(q\), so the factoring
  morphism is identically \(0\), i.e.\ \(\varphi \equiv 0\). Hence \(h = f \circ p + g \circ q\).
  Uniqueness: restricting any such decomposition to the axes recovers \(f\) and \(g\) as the
  axis-restrictions of \(h\).
\end{proof}
% SOURCE QUOTE PROOF: "PROOF. The regular map α will send the k-rational point 0 of A to a
% k-rational point b of B. After composing α with translation by −b, we may assume that
% α(0)=0. Consider the map φ:A×A→B, φ(a,a′)=α(a+a′) − α(a) − α(a′). By this we mean that φ is
% the difference of the two regular maps A×A →(m)→ A →(α)→ B and A×A →(α×α)→ B×B →(m)→ B, which
% is a regular map. Then φ(A×0)=0=φ(0×A) and so φ=0. This means that α is a homomorphism."
% [Milne, Cor 1.2, \S I.1, doc p.~9.]
\begin{proof}
  \leanok
  \textit{Source: Milne, Abelian Varieties, Corollary~1.2.}
  The image \(\alpha(0_A)\) is a \(\bar k\)-point \(b\) of \(B\); composing \(\alpha\) with translation by
  \(-b\) we may assume \(\alpha(0_A) = 0_B\), and it suffices to show such an \(\alpha\) is a
  homomorphism. Form
  \[ \varphi \colon A \times A \to B, \qquad \varphi(a, a') = \alpha(a + a') - \alpha(a) - \alpha(a'), \]
  the difference of the two regular maps \(A \times A \xrightarrow{m} A \xrightarrow{\alpha} B\)
  and \(A \times A \xrightarrow{\alpha \times \alpha} B \times B \xrightarrow{m} B\). Then
  \(\varphi(0_A, 0_A) = 0\) and \(\varphi\) vanishes on both axes,
  \(\varphi(A \times 0) = 0 = \varphi(0 \times A)\) (since \(\alpha(0) = 0\)). By the additive
  decomposition \texttt{lem:hom_additivity_over_product} (with \(V = W = A\)), the axis-restrictions
  of \(\varphi\) being \(0\) force \(\varphi \equiv 0\); equivalently, \(\varphi\) collapses the complete
  axis \(A \times 0\) and vanishes on \(0 \times A\), so Rigidity gives \(\varphi = 0\). Hence
  \(\alpha(a + a') = \alpha(a) + \alpha(a')\), i.e.\ \(\alpha\) is a homomorphism.
\end{proof}
\begin{remark}
  \label{rmk:thm32_codim1_mathlib_gap}
  \textbf{Formalisation risk: the pure-codimension-\(1\) step (Lemma 3.3) has no Mathlib
  Weil-divisor support.} Theorem~3.1's everywhere-extension-outside-codimension-\(2\) is within
  reach of Mathlib's valuative criterion of properness
  (\texttt{ValuativeCriterion}/\texttt{IsProper.of\_valuativeCriterion}), applied at the
  discrete valuation ring \(\mathcal O_{V, Z}\) of each codimension-\(1\) point \(Z\). But Lemma~3.3,
  which rules out codimension-\(1\) indeterminacy for maps into a group variety, is proved in
  Milne via Weil-divisor theory (\(\mathrm{div}(f)_0 - \mathrm{div}(f)_1\), prime divisors,
  tangent cones) for which Mathlib currently has \emph{no} usable API at this generality. This
  is the route's riskiest sub-build. An open question for the prover: whether a \emph{pointwise}
  valuative extension --- showing, at each codimension-\(1\) point \(Z\) separately, that the
  morphism \(\Spec k(V) \to A\) extends over \(\Spec \mathcal O_{V, Z}\) using only that \(A\) is
  proper and a group (so the difference-map argument can be run \emph{locally} at \(Z\) without
  globally building \(\mathrm{div}(f)\)) --- side-steps building divisor theory entirely.
  \textbf{This surface extension is \emph{not} on the genus-\(0\) critical path (corrected, iter-164).}
  An earlier draft routed the genus-\(0\) base case through this extension, by collapsing the
  additive-defect map \(\psi(x, y) = f(x + y) - f(x) - f(y)\) on the \emph{affine} surface
  \(\mathbb G_a \times \mathbb G_a\) and then extending \(\psi\) over the complete surface
  \(\mathbb P^1 \times \mathbb P^1\) --- which genuinely \emph{would} need Theorem~3.2's surface case
  (Milne Theorem~3.4 / Lemma~3.3). That detour is unnecessary. The 4th route
  (\cref{rmk:base_case_fourth_route}) replaces the two-variable affine defect map by the
  \emph{total} translation action \(\sigma \colon \mathbb P^1 \times \mathbb G_a \to \mathbb P^1\)
  (\texttt{def:gaTranslationP1}): \(h := \sigma \fatsemi f\) is a total morphism on \(\mathbb P^1
  \times \mathbb G_a\) (no extension needed), and the complete-source factor is \(\mathbb P^1\)
  itself, so the additive decomposition \texttt{lem:hom_additivity_over_product} applies directly.
  Consequently \cref{thm:rational_map_to_av_extends} (Theorem~3.2), the codim-\(1\) indeterminacy
  sub-build (Lemma~3.3), and the Weil-divisor gap discussed above are all confined to
  \textbf{Route~A} (the positive-genus Albanese universal property, Milne \S III.6); they are
  retained in this chapter solely as Route-A inputs and may be deferred until Route~A consumes
  them. The genus-\(0\) close no longer depends on any of them.
\end{remark}
\section{The genus-\(0\) base objects and the \(\mathbb G_m\)-scaling action (the primary route)}
\begin{definition}
[Concrete genus-\(0\) base objects: \(\mathbb P^1\), \(\mathbb G_a\), \(\mathbb G_m\) over \(\bar k\)]
  \label{def:genus0_base_objects}
  \lean{AlgebraicGeometry.ProjectiveLineBar}
  % Archon-original Lean encoding of standard objects; the mathematics is textbook (Hartshorne
  % I.2--3 for the projective line; the group laws on \(\mathbb A^1\) are elementary). No external
  % % SOURCE QUOTE is needed for the Lean encoding itself.
  Work over \(\Spec \bar k\) with \(\bar k\) algebraically closed. The primary \(\mathbb G_m\)-scaling
  route (and its \(\mathbb G_a\)-additive alternative) needs only three concrete objects of
  \(\mathrm{Over}\,(\Spec \bar k)\):
  \begin{itemize}
    \item \(\mathbb P^1 = \mathbb P^1_{\bar k}\), the projective line --- a smooth proper
      geometrically irreducible curve of genus \(0\) (Mathlib: projective space, or \(\mathrm{Proj}\)
      of the graded polynomial ring; \lean name \texttt{ProjectiveLineBar}). It carries
      the distinguished \(\bar k\)-points \(0\), \(1\), \(\infty\) and the two standard affine charts
      \(\mathbb A^1 = \mathbb P^1 \setminus \{\infty\}\) (coordinate \(x\)) and
      \(\mathbb P^1 \setminus \{0\}\) (coordinate \(u = 1/x\)).
    \item \(\mathbb G_a = \mathbb A^1_{\bar k}\), the affine line as the \emph{additive} group object
      (Mathlib: \texttt{AffineSpace}\,\(\mathbb A(1; \Spec \bar k)\) with the \(\mathtt{GrpObj}\)
      additive structure \((x, y) \mapsto x + y\), identity \(0\); \lean name \texttt{Ga}).
      As a scheme \(\mathbb G_a = \mathbb A^1 = \mathbb P^1 \setminus \{\infty\}\), and the open
      immersion \(\mathbb G_a \hookrightarrow \mathbb P^1\) has dense (cofinite) image.
    \item \(\mathbb G_m = \mathbb A^1 \setminus \{0\} = \mathbb P^1 \setminus \{0, \infty\}\), the
      \emph{multiplicative} group object (group law \((x, y) \mapsto xy\), identity \(1\); \lean name
      \texttt{Gm}). Its open immersion into \(\mathbb P^1\) also has dense image.
  \end{itemize}
  All three are of finite type over \(\bar k\); \(\mathbb P^1\) is proper, while \(\mathbb G_a\) and
  \(\mathbb G_m\) are affine (hence not complete). These are the only ``genus-\(0\)'' geometric inputs
  the base case needs.
\end{definition}
% Per-decl Lean coverage hooks for the components of \cref{def:genus0_base_objects}. The bundled
% definition above pins \texttt{ProjectiveLineBar} in its own \lean{...} block; the sibling
% definition blocks below pin the remaining Lean targets one-by-one so the blueprint-doctor can
% statically link each declaration.
% NOTE (iter-172 encoding clarification): the Lean encoding of \(\mathbb G_m\) realises
% \(\mathbb A^1 \setminus \{0\}\) as the AFFINE scheme \(\Spec(\mathrm{Localization.Away}\,t)\)
% rather than as the basic open subscheme \(D(t) \subset \mathbb A^1\). The two are
% canonically isomorphic; the affine-Spec choice is preferred in this project because
% downstream consumers (the scaling action \(\sigma_\times\) of
% \texttt{def:gaTranslationP1}, the additivity decomposition input
% \(W = \mathbb G_m\) of \texttt{lem:hom_additivity_over_product}) need an affine scheme,
% not just an open subscheme of one.
\begin{lemma}
[\(\mathbb P^1_{\bar k}\) is proper over \(\Spec \bar k\)]
  \label{lem:projectiveLineBar_isProper}
  \lean{AlgebraicGeometry.projectiveLineBar_isProper}
  \uses{def:genus0_base_objects}
  % Archon-original instance; the mathematics is textbook (Hartshorne
  % II.4 + the standard properness of Proj of a finitely generated graded
  % algebra over a Noetherian base; here the base is a field, hence trivially
  % Noetherian). No external SOURCE QUOTE needed.
  The structure morphism \((\mathtt{ProjectiveLineBar}\,\bar k).\mathrm{hom}
  \colon \mathbb P^1_{\bar k} \to \Spec \bar k\) is proper.
  \medskip
  \noindent\textbf{Proof sketch.} Unfold \((\mathtt{ProjectiveLineBar}\,\bar k)
  .\mathrm{hom}\) as the composition
  \(\Proj.\mathrm{toSpecZero}\,(\mathrm{projectiveLineBarGrading}\,\bar k)
  \fatsemi \Spec.\mathrm{map}\,(\mathrm{algebraMap}\,\bar k\,
  \uparrow{\mathcal A_0})\). The first factor is proper by Mathlib's
  \(\mathtt{Proj.instIsProperToSpecZero}\): the polynomial algebra
  \(\bar k[X_0, X_1]\) is finite type over its degree-\(0\) piece
  \(\mathcal A_0\), and the \(\mathtt{IsScalarTower}\) +
  \(\mathtt{Algebra.FiniteType}\) chain (via
  \(\mathtt{Algebra.FiniteType.of\_restrictScalars\_finiteType}\)) supplies the
  finite-type witness. The second factor \(\Spec.\mathrm{map}\,(\mathrm{algebraMap}
  \,\bar k\,\uparrow{\mathcal A_0})\) is an isomorphism: the algebra map
  \(\bar k \to \uparrow{\mathcal A_0}\) is bijective because \(\mathcal A_0\)
  is the constants subalgebra of \(\bar k[X_0, X_1]\) (injectivity via
  \(\mathtt{MvPolynomial.C\_injective}\); surjectivity via
  \(\mathtt{MvPolynomial.homogeneousComponent\_of\_mem}\) +
  \(\mathtt{homogeneousComponent\_zero}\)). Composition of proper with iso is
  proper. LOC: \(\sim 40\) (already landed axiom-clean).
\end{lemma}
\subsection*{Chart cover and chart-ring iso (formalisation infrastructure)}
The next three declarations are iter-168 formalisation infrastructure that the chartwise
construction of the \(\mathbb G_m\)-scaling action \(\sigma_\times\) (\texttt{def:gaTranslationP1})
and its downstream consumers (\texttt{prop:morphism_P1_to_AV_constant}) ultimately rest on.
They realise, at the Lean level, the standard ``affine chart of \(\mathbb P^1\) is \(\mathbb
A^1\)'' picture: a two-chart affine open cover of \(\overline{\mathbb P^1}\) by \(D_+(X_0)\) and
\(D_+(X_1)\); the identification of each chart ring with the one-variable polynomial ring
\(\bar k[u]\); and the resulting reducedness of \(\overline{\mathbb P^1}\) as a scheme. None
of them is on the genus-\(0\) critical path conceptually, but each is consumed by the Lean
formalisation downstream.
\subsection*{Iter-195 analogist recipe --- \(\mathtt{Proj.awayι\_app\_basicOpen}\) port (Lane E close)}
% SOURCE: Mathlib `Mathlib/AlgebraicGeometry/AffineScheme.lean:560-564`
% (the `IsAffineOpen.fromSpec_app_self` formula) and
% `Mathlib/AlgebraicGeometry/ProjectiveSpectrum/Basic.lean:143-185`
% (the `basicOpenToSpec_app_top` + `basicOpenIsoSpec` + `basicOpenIsoAway`
% chain). Read by the iter-195 mathlib-analogist
% `lane-e-proj-appiso-pivot` (cross-domain-inspiration mode);
% verdict ANALOGUE_FOUND.
% NOTE (iter-196 blueprint-writer avr-lane-e-recipe): this subsection records
% the explicit Mathlib-named recipe the iter-196 prover dispatches to close the
% Lane E chart-1 `appTop` residual at line 273 of
% `AlgebraicJacobian/AbelianVarietyRigidity.lean` (the typed `sorry` carried by
% \texttt{kbarChart1Ring\_specMap\_fac}). The iter-196 progress-critic primary
% corrective is precisely \emph{explicit Lean API specification in the
% blueprint}: the analogist file
% \texttt{analogies/lane-e-proj-appiso-pivot.md} is too long and cross-domain
% to feed the prover directly. The two new \(\mathtt{\lean}\) pins below
% (\texttt{lem:awayi_app_basicOpen} and
% \texttt{lem:awayi_appIso_top_inv_apply_isLocElem}) name the new declarations
% the prover will land; the third block records the consumer dispatch at the
% existing typed sorry.
% Lane E history (iter-187 -- iter-195 STUCK). After the iter-187 pivot to
% the separated-locus alternative (\texttt{lem:gmscaling_chart_agreement}, path
% (III.c)), the chart-\(1\) close devolved to a single \(\mathtt{Proj.appIso}\)
% evaluation residual: identify
% \(((\mathrm{Proj.awayι}\,X_1).\mathrm{appIso}\,\top).\mathrm{inv}\,(\mathtt{isLocElem})\)
% with \([X_0/X_1] \in \mathrm{HomogeneousLocalization.Away}\,\mathcal A\,X_1\). The
% iter-188/189/190/191 prover attempts (direct \(\mathtt{simp + rw}\) chains
% through \(\mathtt{Proj.basicOpen\_eq\_basicOpenIso} \blacktriangleright
% \mathtt{Proj.appIso\_inv}\)) hit Lean term-unification timeouts; iter-192
% factored through the named helper \texttt{lem:iotaGm_chart1_appIso_eval} but
% that helper carries the same residual; iter-193 re-routed through
% \(\mathtt{IsOpenImmersion.lift\_uniq}\) (eliminating the outer
% morphism-level reasoning, leaving the inner \(\mathtt{appTop}\) residual
% intact); iter-194 structurally isolated the residual at the consumer
% \(\mathtt{kbarChart1Ring\_specMap\_fac}\) (line 273 of
% \(\mathtt{AlgebraicJacobian/AbelianVarietyRigidity.lean}\)) without
% closing it. The iter-195 \texttt{cross-domain-inspiration} analogist found
% the right Mathlib model: the substantive content of the residual is the
% \(\mathrm{Proj}\)-side mirror of
% \(\mathtt{IsAffineOpen.fromSpec\_app\_self}\)
% (\(\mathtt{Mathlib/AlgebraicGeometry/AffineScheme.lean:560\text{--}564}\)),
% whose proof shape transports verbatim. The three blocks below capture the
% port and its consumer dispatch.
% NOTE (iter-197 blueprint-writer): the iter-196 prover identified the
% intermediate \(\mathrm{Proj.basicOpenIsoSpec}.\mathrm{inv}.\mathrm{app}\,\top\)
% identity as the missing named helper whose absence caused the dependent-motive
% blocker on \texttt{lem:awayi_app_basicOpen}. It is now pinned below as
% \texttt{lem:basicOpenIsoSpec_inv_app_top}, immediately preceding the section-level
% port. The two iter-196 landed primitives (the preimage identity and the
% \(\mathrm{Spec.map} \fatsemi \mathrm{fromSpec}\) bridge) are also pinned just
% above for blueprint visibility.
\noindent\textbf{Consumer dispatch (iter-196 prover).} The existing
typed sorry at line \(273\) of
\(\mathtt{AlgebraicJacobian/AbelianVarietyRigidity.lean}\) in
\(\mathtt{kbarChart1Ring\_specMap\_fac}\) reduces, after
\texttt{lem:awayi_appIso_top_inv_apply_isLocElem} is in scope, to a
forward equation between two ring-map expressions on
\(\mathtt{appTop}\). The iter-196 prover dispatch is:
\begin{enumerate}
  \item Rewrite the LHS using
    \(\mathtt{simp\ only}\,[\,\mathtt{Proj.awayι\_appIso\_top\_inv\_apply\_isLocElem}\,]\)
    to substitute the closed form supplied by
    \texttt{lem:awayi_appIso_top_inv_apply_isLocElem}.
  \item Identify the RHS by \(\mathtt{ext\_of\_isAffine}\) plus
    \(\mathtt{awayToSection\_apply}\), exhibiting the resulting ring map as
    \(\mathrm{IsLocalization.map}\,\_\,\mathrm{HomogeneousLocalization.val}\)
    on the canonical generator.
  \item Close by definitional unfolding of \(\mathtt{kbarChart1Ring}\)
    and the canonical
    \(\mathrm{MvPolynomial.eval}\,(\lambda \_ \mapsto 1)\) collapse,
    matching the chart-\(1\) ring map.
\end{enumerate}
This dispatch is the substantive content the iter-188 -- iter-194
attempts converged on; it is now packaged in the two named helpers
\texttt{lem:awayi_app_basicOpen} and
\texttt{lem:awayi_appIso_top_inv_apply_isLocElem}, with the consumer-site
discharge estimated at \(\sim 5\)--\(10\) LOC. The same recipe also closes
the residual sorry inside \texttt{lem:iotaGm_chart1_appIso_eval} (line
\(481\) of \(\mathtt{AlgebraicJacobian/AbelianVarietyRigidity.lean}\)),
which shares the same \(\mathtt{Proj.appIso}\) evaluation residual modulo
the iso-chain functoriality already discharged in stages 5--8 of that
helper's proof.
\subsection*{Scaffolds for the body of \(\sigma_\times = \mathtt{gmScalingP1}\) (iter-171 chart-glue split)}
% NOTE (iter-173, blueprint-writer g0bo-pin-scaffolds): the iter-171 option-(c) inline chart-glue
% commitment factored the body of \texttt{gmScalingP1} (\texttt{def:gaTranslationP1}) into a
% \texttt{Scheme.Cover.glueMorphisms} over the pullback cover, with three top-level named
% scaffold sorries naming the per-chart scheme morphism, its cocycle on overlap, and the
% \texttt{Over}-side structure-morphism coherence. The four blocks below pin those scaffolds
% one-by-one so the dependency graph and \texttt{sync\_leanok} can track each declaration
% independently (the iter-172 lean-vs-blueprint-checker \texttt{g0bo172} MAJOR finding).
% Iter-173 prover work on the body of \(\sigma_\times\) targets exactly these three obligations.
% NOTE (iter-186 writer, gmscaling-chart-agreement-expansion): pinning the two project-side
% simp helpers `pullback_map_{fst,snd}_proj` (analogist Recipe 1, iter-184) here so that the
% expanded proof sketch of \texttt{lem:gmscaling_chart_agreement} below can `\texttt{...}` them
% without an unresolved cross-reference.
  
% NOTE (iter-186 writer, gmscaling-chart-agreement-expansion): mini-blocks below pin the
% four product-stability instances on \(\mathbb P^1 \otimes \mathbb G_m\) that
% \texttt{prop:morphism_P1_to_AV_constant}'s Lean basepoint helper consumes (one local-finiteness
% instance, one reducedness instance, and the two geometric-irreducibility instances). They
% are exported from `AlgebraicJacobian/Genus0BaseObjects/GmScaling.lean`. The two
% sorry-bearing instances (`gm_geomIrred`, `projGm_isReduced`) carry the Mathlib gap recorded
% in \texttt{lem:gmscaling_chart_agreement}'s iter-186 NOTE.
% SOURCE QUOTE PROOF: "PROOF. According to (3.2), α extends to a regular map on the whole of A¹.
% After composing α with a translation, we may suppose that α(0)=0. Then α is a homomorphism,
% α(x+y)=α(x)+α(y); all x,y ∈ A¹(k^al)=k^al. But A¹−{0} is also a group variety, and similarly,
% α(xy)=α(x)+α(y)+c; all x,y ∈ A¹(k^al)=k^al. This is absurd, unless α is constant." [Milne,
% Prop 3.9, \S I.3, doc p.~19--20.]
\begin{proof}
  \textit{Source: Milne, Abelian Varieties, Proposition~3.9.}
  Write \(A\) additively (an abelian variety is commutative). We give the conceptual
  \emph{completeness} proof as the abstract argument, then record Milne's elementary \(\mathbb
  G_a/\mathbb G_m\) incompatibility, which applies whenever \(\varphi\) extends to a morphism on
  \(\mathbb P^1\) (as it does in the genus-\(0\) application).
  \textbf{Completeness proof.} The image \(\varphi(\mathbb G_a)\) is a connected subgroup of \(A\)
  (image of the connected \(\mathbb G_a\) under a homomorphism). As the image of a homomorphism of
  algebraic groups it is a \emph{closed} subgroup of \(A\), hence \emph{complete} (closed in the
  complete \(A\)); as a homomorphic image of the \emph{affine} \(\mathbb G_a\) it is moreover
  \emph{affine}. A connected variety that is simultaneously complete and affine over \(\bar k\) is a
  single point: its ring of global regular functions is a field finite over \(\bar k\), hence equal
  to \(\bar k\) (the \(\Gamma = \bar k\) collapse already packaged for proper integral sources in
  \cref{lem:eq_comp_of_isAffine_of_properIntegral}). Therefore \(\varphi(\mathbb G_a)\) is one point;
  being a homomorphism, \(\varphi\) sends \(0\) there, so the point is \(0_A\) and
  \(\varphi = \mathtt{toUnit}\,\mathbb G_a \fatsemi \eta_A\).
  \medskip
  \noindent\textbf{Milne's elementary \(\mathbb G_a/\mathbb G_m\) incompatibility (when \(\varphi\)
  extends to \(\mathbb P^1\)).} Suppose \(\varphi = f|_{\mathbb G_a}\) for a morphism
  \(f \colon \mathbb P^1 \to A\) (the genus-\(0\) situation). Set \(g := f - f(1)\), so \(g(1) = 0_A\), and
  form \(h_\times := \sigma_\times \fatsemi g \colon \mathbb P^1 \times \mathbb G_m \to A\) with the
  scaling companion \(\sigma_\times\) of \texttt{def:gaTranslationP1}. Applying the additive
  decomposition \texttt{lem:hom_additivity_over_product} with \(V = \mathbb P^1\) (complete),
  \(W = \mathbb G_m\), base points \(1 \in \mathbb P^1\) and \(1 \in \mathbb G_m\) (so
  \(h_\times(1, 1) = g(1) = 0_A\)), the axis restrictions are \(x \mapsto g(x)\) and
  \(\lambda \mapsto g(\lambda)\), whence \(g(\lambda x) = g(x) + g(\lambda)\) on \(\bar k^\times\);
  re-expanding \(g\) gives the multiplicative relation
  \(\varphi(xy) = \varphi(x) + \varphi(y) + c\) with \(c = -\varphi(1)\). Now \(\varphi\) is at once
  additive, \(\varphi(x + y) = \varphi(x) + \varphi(y)\), and ``multiplicative-additive'',
  \(\varphi(xy) = \varphi(x) + \varphi(y) + c\). Computing \(\varphi((x + x')\,y)\) in two ways for
  \(x, x' \in \bar k\), \(y \in \bar k^\times\),
  \[ \varphi((x+x')y) = \varphi(xy) + \varphi(x'y) = \varphi(x) + \varphi(x') + 2\varphi(y) + 2c, \]
  while also \(\varphi((x+x')y) = \varphi(x+x') + \varphi(y) + c = \varphi(x) + \varphi(x') +
  \varphi(y) + c\); equating gives \(\varphi(y) = -c = \varphi(1)\) for every \(y \in \bar k^\times\).
  Hence \(\varphi\) is constant on the dense \(\mathbb G_m \subset \mathbb G_a\), so constant on
  \(\mathbb G_a\) (density, \(A\) separated), with value \(\varphi(0) = 0_A\). This is Milne's argument
  verbatim and uses only the proven \texttt{lem:hom_additivity_over_product} plus the scaling action,
  avoiding the ``image is a closed subgroup'' input of the completeness proof. \qedhere
\end{proof}
\begin{remark}
  \label{rmk:base_case_fourth_route}
  \textbf{The primary route is the \(\mathbb G_m\)-scaling shortcut; why no Theorem~3.2 and no
  \(\mathrm{Hom}(\mathbb G_a, A) = 0\) are needed.} The genus-\(0\) base case ``\(f \colon \mathbb P^1
  \to A\) is constant'' is closed \emph{directly} by the proven additive decomposition
  \texttt{lem:hom_additivity_over_product} composed with a total \(\mathbb G\)-action on \(\mathbb P^1\),
  with \emph{no} appeal to Milne's Theorem~3.2 / Lemma~3.3 (rational-map extension), the theorem of
  the cube, Auslander--Buchsbaum, or differentials. The reason the earlier draft thought
  Theorem~3.2 was needed --- that the additive defect \(\psi(x,y) = f(x+y) - f(x) - f(y)\) lives only
  on the affine \(\mathbb G_a \times \mathbb G_a\) and must be extended over \(\mathbb P^1 \times
  \mathbb P^1\) --- evaporates once one composes with a \emph{total} action morphism of
  \texttt{def:gaTranslationP1}: the composite is already a total morphism out of \(\mathbb P^1 \times
  \mathbb G\), and the complete factor is \(\mathbb P^1\), exactly what
  \texttt{lem:hom_additivity_over_product} (whose Lean signature needs only the first factor
  complete) consumes.
  \medskip
  \noindent\textbf{Primary route --- the scaling shortcut (strategy decision of iter-164).} The
  \emph{multiplicative} scaling action \(\sigma_\times \colon \mathbb P^1 \times \mathbb G_m \to
  \mathbb P^1\), \((x, \lambda) \mapsto \lambda x\) (\texttt{def:gaTranslationP1}), closes the base case
  in one shot, \emph{without} the additive homomorphism step or \texttt{lem:hom_Ga_to_av_trivial} at
  all. Normalise \(f(0) = 0_A\) (translate; here \(0\) is a scaling fixed point). Form \(h_\times :=
  \sigma_\times \fatsemi f\) and apply \texttt{lem:hom_additivity_over_product} with \(V = \mathbb P^1\),
  \(W = \mathbb G_m\), base points \(0 \in \mathbb P^1\) and \(1 \in \mathbb G_m\) (so \(h_\times(0, 1) =
  f(0) = 0_A\)). Because \(0\) is a \emph{fixed point} of scaling, the \(W\)-axis restriction collapses:
  \(\lambda \mapsto h_\times(0, \lambda) = f(\lambda \cdot 0) = f(0) = 0_A\). The decomposition
  therefore reads \(h_\times = p \fatsemi f\) with the second summand the group identity, i.e.
  \[ \sigma_\times \fatsemi f = \mathrm{pr}_1 \fatsemi f \colon \mathbb P^1 \times \mathbb G_m \to A,
     \qquad\text{``}f(\lambda x) = f(x)\text{''}. \]
  Precomposing with \(\lambda \mapsto (1, \lambda)\) gives \(f|_{\mathbb G_m} = \text{const}\,f(1)\);
  since \(\mathbb G_m\) is dense in \(\mathbb P^1\) and \(A\) is separated, \(f\) is constant on \(\mathbb
  P^1\) (the dominant-source/separated-target handle \texttt{thm:GrpObj_eq_of_eqOnOpen}, i.e.\ the
  \(\mathtt{ext\_of\_isDominant}\) family already used by \texttt{rigidity\_core}). This uses
  \emph{only} the proven \texttt{lem:hom_additivity_over_product}, the total morphism \(\sigma_\times\),
  and density --- no \(\mathbb G_a\) additivity, no \texttt{lem:hom_Ga_to_av_trivial}, and in
  particular \emph{no} ``image of a group homomorphism is a closed subgroup'' input (a deep Mathlib
  gap). It is the least-Mathlib-blocked realisation of the base case and is the route the live
  consumer \texttt{prop:morphism_P1_to_AV_constant} formalises.
  \medskip
  \noindent\textbf{Retained alternative --- the \(\mathbb G_a\)-additive route.} The Milne-faithful
  argument routes through the translation action \(\sigma \colon \mathbb P^1 \times \mathbb G_a \to
  \mathbb P^1\): \(\sigma \fatsemi f\) is total, the additive decomposition makes \(f|_{\mathbb G_a}\) a
  
  homomorphism \(\mathbb G_a \to A\) (\texttt{lem:hom_from_Ga_trivial}), and a homomorphism from the
  
  additive group into an abelian variety is trivial (\texttt{lem:hom_Ga_to_av_trivial}, Milne's
  \(\mathbb G_a/\mathbb G_m\) incompatibility, Proposition~3.9). This reaches the same conclusion but
  needs the triviality lemma; it is retained for fidelity to Milne and is \emph{off} the genus-\(0\)
  critical path.
\end{remark}
% SOURCE QUOTE PROOF: "PROOF. According to (3.2), α extends to a regular map on the whole of A¹.
% After composing α with a translation, we may suppose that α(0)=0. Then α is a homomorphism,
% α(x+y)=α(x)+α(y); all x,y ∈ A¹(k^al)=k^al. But A¹−{0} is also a group variety, and similarly,
% α(xy)=α(x)+α(y)+c; all x,y ∈ A¹(k^al)=k^al. This is absurd, unless α is constant." [Milne,
% Prop 3.9, \S I.3, doc p.~19--20.]
\begin{proof}
  \textit{Source: Milne, Abelian Varieties, Proposition~3.9.}
  Normalise \(f(0) = 0_A\) by a translation in \(A\) (un-translate at the end). The proof is the
  \emph{4th route} (\cref{rmk:base_case_fourth_route}): the additive defect map is never built on
  the affine surface and never extended; instead we precompose \(f\) with the \emph{total}
  translation action \(\sigma\), and the complete factor is \(\mathbb P^1\) itself, so the proven
  additive decomposition applies directly. Milne's Theorem~3.2
  (\cref{thm:rational_map_to_av_extends}) is \emph{not} used.
  \textbf{\(f|_{\mathbb G_a}\) is an additive homomorphism, via the total action \(\sigma\).} Form
  \[ h := \sigma \fatsemi f \colon \mathbb P^1 \times \mathbb G_a \to A, \qquad
     \text{``}h(x, y) = f(x + y)\text{''}, \]
  using the translation action \(\sigma\) of \texttt{def:gaTranslationP1}; this is a \emph{total}
  morphism (no rational-map extension needed). Apply the additive decomposition
  \texttt{lem:hom_additivity_over_product} (Milne Cor~1.5) --- whose Lean signature requires only the
  \emph{first} factor complete --- with \(V = \mathbb P^1\) (proper), \(W = \mathbb G_a\), and base
  points \(0 \in \mathbb P^1\), \(0 \in \mathbb G_a\), noting \(h(0, 0) = f(0) = 0_A\). The axis
  restrictions are \(x \mapsto h(x, 0) = f(x)\) and \(y \mapsto h(0, y) = f(y)\), so the decomposition
  \(h = (p \fatsemi f) + (q \fatsemi (f|_{\mathbb G_a}))\) reads, on \(\bar k\)-points,
  \[ f(x + y) = f(x) + f(y), \qquad x \in \mathbb P^1(\bar k),\ y \in \mathbb G_a(\bar k) = \bar k. \]
  Restricting \(x\) to \(\mathbb G_a\) shows \(f|_{\mathbb G_a} \colon \mathbb G_a \to A\) is a
  homomorphism of group objects (it sends \(0 \mapsto 0_A\) and respects addition).
  \textbf{A homomorphism \(\mathbb G_a \to A\) is trivial.} By \texttt{lem:hom_Ga_to_av_trivial}
  (the \(\mathrm{Hom}(\mathbb G_a, A) = 0\) lemma, the content of Milne's \(\mathbb G_a/\mathbb G_m\)
  incompatibility), \(f|_{\mathbb G_a}\) is the constant morphism at \(0_A\). Hence \(f\) is constant on
  the dense open \(\mathbb G_a \subset \mathbb P^1\); since \(\mathbb G_a\) is dense and \(A\) is
  separated, \(f\) is constant on all of \(\mathbb P^1\), with value \(f(0) = 0_A\) (i.e.\ the original
  base value before the normalising translation). \qedhere
  \medskip
  \noindent\emph{Note: this route is the demoted alternative.} The primary base-case proof is the
  \(\mathbb G_m\)-scaling shortcut now installed directly as the proof of
  \texttt{prop:morphism_P1_to_AV_constant} (\cref{rmk:base_case_fourth_route}): the multiplicative
  scaling action \(\sigma_\times\) of \texttt{def:gaTranslationP1} gives \(\sigma_\times \fatsemi f =
  \mathrm{pr}_1 \fatsemi f\) (the \(W\)-axis collapses because \(0\) is a scaling fixed point), so
  
  \(f|_{\mathbb G_m}\) is constant and \(f\) is constant by density --- bypassing both the \(\mathbb
  G_a\)-additivity step and \texttt{lem:hom_Ga_to_av_trivial}, and avoiding the ``image is a closed
  subgroup'' Mathlib gap. The present \(\mathbb G_a\)-additive lemma is retained only for fidelity to
  \uses{lem:hom_additivity_over_product, def:gaTranslationP1, lem:gmScaling_fixes_zero}
  Milne; it is not on the genus-\(0\) critical path.
\end{proof}
\section{Morphisms from \(\mathbb P^1\) to an abelian variety are constant}
% SOURCE QUOTE PROOF: "PROOF. We may suppose that \(k\) is algebraically closed. By assumption
% there is a rational map \(\mathbb A^n \dashrightarrow V\) with dense image, and the composite
% of this with \(\alpha\) extends to a morphism $\beta \colon \mathbb P^1 \times \cdots \times
% \mathbb P^1 \to A$. An induction argument, starting from Corollary 1.5, shows that there are
% regular maps \(\beta_i \colon \mathbb P^1 \to A\) such that $\beta(x_1, \ldots, x_d) = \sum
% \beta_i(x_i)\(, and the lemma shows that each \)\beta_i$ is constant."
\begin{proof}
  This is the single-factor (\(d = 1\)) instance of Milne's Proposition~3.10, where the
  unirational variety is \(\mathbb P^1\) itself. We give the project's \textbf{primary} proof, the
  \emph{\(\mathbb G_m\)-scaling shortcut}: it proves constancy \emph{in one shot} from the proven
  additive decomposition (\texttt{lem:hom_additivity_over_product}) and a scaling fixed point,
  without the intermediate ``\(f|_{\mathbb G_a}\) additive'' step and \emph{without} the
  \(\mathrm{Hom}(\mathbb G_a, A) = 0\) triviality lemma. The theorem of the cube is \emph{not} used
  (\cref{rmk:cube_not_needed}), and the map \(f\) is already a morphism out of the complete
  nonsingular \(\mathbb P^1\), so Milne's Theorem~3.2 (\cref{thm:rational_map_to_av_extends}) is
  \emph{not} invoked either.
  \textbf{Reduction to a base point.} Choose the scaling fixed point \(0 \in \mathbb G_a \subset
  \mathbb P^1\) (one of the two fixed points of \(\sigma_\times\)). Translating in the group \(A\) by
  \(-f(0)\), we may assume \(f(0) = 0_A\); un-translating at the end restores the original value.
  \textbf{The scaling shortcut.} Form the \emph{total} morphism
  \[ h := \sigma_\times \fatsemi f \colon \mathbb P^1 \times \mathbb G_m \to A, \qquad
     \text{``}h(x, \lambda) = f(\lambda x)\text{''}, \]
  using the multiplicative scaling action \(\sigma_\times\) of \texttt{def:gaTranslationP1} (a genuine
  total morphism --- no rational-map extension needed). Apply the additive decomposition
  \texttt{lem:hom_additivity_over_product} (Milne Cor~1.5) --- whose Lean signature requires only the
  \emph{first} factor complete --- with \(V = \mathbb P^1\) (proper), \(W = \mathbb G_m\), and base
  points \(0 \in \mathbb P^1\), \(1 \in \mathbb G_m\); here \(h(0, 1) = f(1 \cdot 0) = f(0) = 0_A\), as
  required. The decomposition writes \(h = (p \fatsemi f_0) + (q \fatsemi g)\) with \(f_0\) the
  \(V\)-axis restriction \(x \mapsto h(x, 1) = f(1 \cdot x) = f(x)\) (so \(f_0 = f\)) and \(g\) the
  \(W\)-axis restriction \(\lambda \mapsto h(0, \lambda) = f(\lambda \cdot 0) = f(0) = 0_A\).
  \textbf{The \(W\)-axis collapses (the load-bearing scaling fixed point).} Because \(0\) is a
  \emph{fixed point} of scaling, \(\sigma_\times(0, \lambda) = 0\) for all \(\lambda\)
  (\texttt{lem:gmScaling_fixes_zero}), so the \(W\)-axis
  restriction \(g\) is the \emph{constant} morphism at \(0_A\) --- the identity of the hom-group
  \(\mathrm{Hom}(\mathbb G_m, A)\). The corresponding summand \(q \fatsemi g\) is therefore the group
  identity, and the decomposition collapses to
  \[ h = p \fatsemi f, \qquad\text{i.e.}\qquad
     \sigma_\times \fatsemi f = \mathrm{pr}_1 \fatsemi f, \qquad
     \text{``}f(\lambda x) = f(x)\text{''} \quad (x \in \mathbb P^1,\ \lambda \in \mathbb G_m). \]
  \textbf{Constancy via density.} Specialise \(x = 1\) (precompose with \(\lambda \mapsto (1,
  \lambda)\)): \(f(\lambda) = f(1)\) for all \(\lambda \in \mathbb G_m\), i.e.\ \(f|_{\mathbb G_m}\) is
  the constant morphism at \(f(1)\). Since \(\mathbb G_m\) is dense in the irreducible \(\mathbb P^1\)
  and \(A\) is separated, the two morphisms \(f\) and \(\mathrm{const}_{f(1)}\), agreeing on the
  non-empty open \(\mathbb G_m\), agree everywhere by the dominant-source/separated-target rigidity
  handle \texttt{thm:GrpObj_eq_of_eqOnOpen} (the same \texttt{ext\_of\_isDominant} family
  \texttt{rigidity\_core} uses). Hence \(f\) is constant, with value \(f(1)\); in the normalised
  picture \(f(0) = 0_A\), so the constant is \(0_A\) (the original \(f(0)\) before translation).
  \medskip
  \noindent\emph{Multi-factor remark.} For the general unirational statement of Proposition~3.10,
  one writes \(\beta(x_1, \ldots, x_d) = \sum_i \beta_i(x_i)\) via the additive decomposition
  (\texttt{lem:hom_additivity_over_product}) and applies this base case to each \(\beta_i \colon
  \mathbb P^1 \to A\); we need only the \(d = 1\) case here. The Milne-faithful \(\mathbb
  G_a\)-additive route (\texttt{lem:hom_from_Ga_trivial}, via \(\mathrm{Hom}(\mathbb G_a, A) = 0\),
  \texttt{lem:hom_Ga_to_av_trivial}) reaches the same conclusion and is retained as an alternative,
  but is \emph{not} on this critical path.
  % NOTE (iter-166 review): The Lean factors the proof into the outer translation reduction
  % (`AlgebraicGeometry.morphism_P1_to_grpScheme_const`, public) plus a private basepoint helper
  % (`AlgebraicGeometry.morphism_P1_to_grpScheme_const_aux`) carrying the body of the
  % \(\mathbb G_m\)-scaling shortcut. The chapter presents the proof at a single granularity; the
  % Lean split is structural-only and preserves the chapter's logical content. The helper carries
  % five honest scaffold sorries (three \(\mathbb P^1 \otimes \mathbb G_m\) product-instance
  % obligations [GeometricallyIrreducible / LocallyOfFiniteType / IsReduced], one
  % \(\mathtt{IsReduced}\ (\mathbb P^1)_{\text{left}}\) obligation, and one
  % \(\mathtt{IsDominant}\ \iota_{\mathbb G_m}.\text{left}\) obligation — the latter strictly
  % downstream of \(\sigma_\times\) landing a body); these are Mathlib bridge work and a Lane-2
  % dependency, not laundered hypotheses. Both `lean-vs-blueprint-checker-avr-iter166` and
  % `lean-auditor-iter166` confirmed: no iter-157-style laundering, basepoint hypothesis \(hf_0\)
  % is genuinely consumed (L975 of the Lean: `rw [← Category.assoc, hcollapse]; exact hf0`),
  % `lean_verify` axioms $\{\mathtt{propext}, \mathtt{sorryAx}, \mathtt{Classical.choice},
  % \mathtt{Quot.sound}\}$ propagate cleanly. Lifts to axiom-clean once the 5 helper sorries
  % discharge.
  % NOTE (iter-167 review): the 5 helper sorries above have been reduced to ONE. The four
  % product / \(\mathbb P^1\) instance obligations are now discharged via Lane A's named exports
  % (`projGm_geomIrred`, `projGm_locallyOfFiniteType`, `projGm_isReduced`,
  % `projectiveLineBar_isReduced` in `AlgebraicJacobian.Genus0BaseObjects`) and resolve by
  % `infer_instance` inside the helper. The remaining $\mathtt{IsDominant}\ \iota_{\mathbb G_m}
  % .\text{left}$ obligation has been promoted to a named file-local lemma
  % `AlgebraicGeometry.iotaGm_isDominant` (still `sorry`, strictly gated on Lane A shipping
  % the chartwise body of \(\sigma_\times = \mathtt{gmScalingP1}\)). So the helper now contains
  % zero inline `sorry`s and the AVR-side count for this proof is exactly one named bridge
  % (`iotaGm_isDominant`). `lean-vs-blueprint-checker-avr-iter167` confirms no laundering and
  % zero excuse-comments anywhere in AVR.lean.
  % NOTE (iter-168 review): of the 4 Lane A product/\(\mathbb P^1\) exports cited by the
  % iter-167 NOTE, `projectiveLineBar_isReduced` is now genuinely proven axiom-clean
  % (iter-168 substantive closure in `AlgebraicJacobian.Genus0BaseObjects`); the prior
  % "Mathlib gap" framing was incorrect — the bridge
  % `HomogeneousLocalization.Away` \(\to\) `IsDomain` is \(<10\) LOC via
  % `Function.Injective.isDomain` on the algebraMap into `Localization.Away`. The other 3
  % cited Lane A exports (`projGm_geomIrred`, `projGm_isReduced`, and the AVR-side
  % NOTE (iter-185 review): per `lean-vs-blueprint-checker-iter185-gmscaling`,
  % `gm_geomIrred` and `projGm_isReduced` are CONFIRMED Mathlib gaps at b80f227:
  % `Algebra.TensorProduct.isDomain_of_isAlgClosed_left` (or any "factor is alg-closed
  % + other is domain ⟹ tensor is domain" bridge) is absent. Both instances stay as
  % typed sorries until upstream lands the bridge, OR a project-side `~30-50` LOC shim
  % via `Function.Injective.isDomain` on the `algebraMap (kbar ⊗ A) (kbar ⊗ A)` chain
  % is contributed. The Lean docstring on each typed-sorry instance documents this gap
  % verbatim; no \(\mathtt{TEMP}\) axiom-by-axiom workaround is being attempted.
  
  % iter-168 also landed two upstream infrastructure pieces in `Genus0BaseObjects.lean` —
  % `projectiveLineBarAffineCover` (the 2-chart affine cover) and the chart-ring iso scaffold
  
  % `homogeneousLocalizationAwayIso` (forward direction axiom-clean; `aux_left` round-trip
  % still `sorry`) — that prepare the chartwise glue for the iter-169 attack on
  % `gmScalingP1`'s body via either the direct `Proj.fromOfGlobalSections` route or the iso
  % route. These new infrastructure pieces are not yet pinned by `\lean{...}` blocks in this
  % chapter; this is a blueprint-writer-iter169 dispatch.
\end{proof}
\section{Genus-\(0\) base-object helper declarations (1-to-1 Lean coverage)}
\label{sec:genus0_helpers}
This section records, for 1-to-1 Lean\(\leftrightarrow\)blueprint coverage, the supporting
declarations of the four \texttt{Genus0BaseObjects} sub-files (\texttt{BareScheme},
\texttt{ChartIso}, \texttt{Points}, \texttt{GmScaling}). These are the construction
helpers, rewrite lemmas, and typeclass instances backing the headline objects \(\mathbb
P^1_{\bar k}\), \(\mathbb G_a\), \(\mathbb G_m\), the chart-ring iso, and the \(\mathbb
G_m\)-scaling action already pinned above. The substantive ones carry a short proof
sketch; the pure plumbing (instances, projection/round-trip components, \(\mathtt{eval}\)
rewrites) carries a one-line justification, its purpose being to carry the dependency edge.
All are Archon-original or directly Mathlib-backed; none quotes external source text.

\subsection*{\texttt{BareScheme.lean}: the bare scheme \(\mathbb P^1\) and its presentation supplement}


\subsection*{\texttt{ChartIso.lean}: the chart-ring iso \(\mathrm{Away}\,\mathcal A\,(X_i) \cong \bar k[u]\)}


\subsection*{\texttt{Points.lean}: standard \(\bar k\)-points and the groups \(\mathbb G_a\), \(\mathbb G_m\)}


\subsection*{\texttt{GmScaling.lean}: chart bridge, cover isos, and the \(\sigma_\times\) cocycle}

